% ----------------------------------------------------------
% Apêndice A — Dataset de Avaliação
% As listagens (apendice_dataset_P/N/gerado.tex) e as tabelas em tabelas/
% são geradas por `pfc-dataset --apendice` e `pfc-auditar --paper`: não editar à mão.
% ----------------------------------------------------------
\chapter{Dataset de Avaliação}
\label{apendice:dataset}

Este apêndice documenta o \textit{dataset} de avaliação da Seção~\ref{sec:dataset}, com o manual de anotação que define o gabarito, as 310 consultas com as leituras aceitas, o mapeamento das consultas aos requisitos, o registro da auditoria e as tabelas de estrutura, composição e concordância dos conjuntos. O arquivo \texttt{data/dataset.json} do repositório (Apêndice~\ref{apendice:codigo}) guarda as mesmas informações em formato legível por máquina. As listagens deste apêndice são geradas das mesmas fontes que produzem esse arquivo (\texttt{manual\_cases.py} e o gerador da camada G), sem transcrição manual.

\paragraph{Como ler as entradas.} Cada consulta traz o identificador, as categorias (Seção~\ref{sec:dataset}), o texto exatamente como é enviado ao modelo e as \textbf{leituras aceitas} do gabarito, na forma canônica da ferramenta. A notação segue o princípio P3 do manual (Seção~\ref{apx:manual}). Nela, ``\texttt{A} ou \texttt{B}'' indica valor alternativo, e ``(opcional)'' indica campo que pode faltar sem erro, mas que, se presente, precisa ter o valor indicado. A forma ``aceita-se também \texttt{city} no lugar de \texttt{state}'' indica leitura estrutural alternativa. As expressões de tempo relativo aparecem pelo nome da regra (``ano corrente'', ``semana passada'') e são resolvidas na data de referência da execução, segundo o Quadro~\ref{quadro:tempo-relativo}. Na camada P, uma nota registra a linha de origem no arquivo do protótipo. Nas camadas P e N, a nota registra também, quando há, a justificativa das leituras alternativas.

\paragraph{Escopo do gabarito.} O gabarito descreve os \textbf{parâmetros} que a consulta deve produzir. A recuperação de produtos fica fora do escopo do trabalho (Seção~\ref{sec:exclusoes-avaliacao}).

% \dtcase / \dtcasenota: bloco inline (não é float) para cada entrada do dataset.
% Não usa \begin{quadro}[htbp!] porque isso exigiria o pacote `float`, que
% conflita com o \newfloat do memoir/abntex2.
\newcommand{\dtcase}[4]{%
  \par\medskip\noindent
  {\small
  \begin{tabular}{|p{2.3cm}|p{12cm}|}
  \hline
  \textbf{ID} & \texttt{#1} \\ \hline
  \textbf{Categorias} & #2 \\ \hline
  \textbf{Consulta} & \textit{``#3''} \\ \hline
  \textbf{Leituras aceitas} & \begin{minipage}[t]{12cm}\vspace{0.05cm}\texttt{#4}\vspace{0.1cm}\end{minipage} \\ \hline
  \end{tabular}}
  \par\medskip
}
\newcommand{\dtcasenota}[5]{%
  \par\medskip\noindent
  {\small
  \begin{tabular}{|p{2.3cm}|p{12cm}|}
  \hline
  \textbf{ID} & \texttt{#1} \\ \hline
  \textbf{Categorias} & #2 \\ \hline
  \textbf{Consulta} & \textit{``#3''} \\ \hline
  \textbf{Leituras aceitas} & \begin{minipage}[t]{12cm}\vspace{0.05cm}\texttt{#4}\vspace{0.1cm}\end{minipage} \\ \hline
  \textbf{Nota} & #5 \\ \hline
  \end{tabular}}
  \par\medskip
}

% ----------------------------------------------------------
% Manual de anotação do dataset (incluído no Apêndice A)
% Fonte: pfc_busca/docs/manual_de_anotacao.md
% ----------------------------------------------------------
\section{Manual de anotação}
\label{apx:manual}

O gabarito de cada consulta foi construído segundo as regras abaixo. Elas são a referência para a redação e a revisão das consultas, para a auditoria automática e para a anotação independente descritas na Seção~\ref{apx:auditoria}; qualquer anotação que as contrarie é considerada erro do gabarito, e não do modelo.

\subsection*{Princípios}

\begin{description}
  \item[P1 --- Evidência explícita.] O gabarito contém somente parâmetros para os quais existe um trecho da consulta que os justifica. Nada é inferido além do que o texto diz (``cartas de Manaus'' não implica \texttt{state = "Amazonas"}).
  \item[P2 --- Forma canônica.] Todo valor é escrito na forma canônica da ferramenta \texttt{buscar\_catalogo}: enumerados exatamente como no \textit{schema}; estados e municípios por extenso, com acentos; datas em ISO 8601.
  \item[P3 --- Leituras múltiplas explícitas.] Quando a consulta admite mais de uma leitura razoável e nem o texto nem a descrição da ferramenta desempatam, o gabarito lista todas as leituras aceitas, sendo a primeira a preferencial. Há três mecanismos: \emph{valor alternativo} (o campo aceita mais de um valor), \emph{campo opcional} (pode faltar; se aparecer, precisa ter o valor indicado) e \emph{leitura estrutural alternativa} (o mesmo trecho pode ocupar campos diferentes, como estado ou município).
  \item[P4 --- Fora do domínio.] Consultas sem intenção de busca no acervo cartográfico brasileiro têm como gabarito ``não chamar a ferramenta''.
  \item[P5 --- Consistência.] A mesma expressão recebe a mesma anotação em todo o \textit{dataset}.
  \item[P6 --- Casos observacionais.] Consultas cujo comportamento correto não é determinável a partir do texto (valor impossível de representar, comportamento esperado discutível) são executadas e reportadas, mas ficam fora das métricas principais. Consultas fora do domínio (P4) não são observacionais, pois o comportamento correto é determinado.
\end{description}

\subsection*{Regras por campo}

\begin{description}
  \item[\texttt{keyword}] Código MI com seus sufixos, sem o prefixo ``MI''; código INOM na grafia com hífens (grafias compactadas, como ``sf22yd'', são normalizadas); nome próprio de carta quando precedido de ``carta'' ou ``folha''. Se o nome da carta coincide com um município, aceita-se também \texttt{city}. Havendo dois códigos, aceita-se qualquer um deles, pois o \textit{schema} comporta um só.
  \item[\texttt{scale}] Escala explícita em qualquer grafia (``1:25.000'', ``1:25000'', ``25k'', ``25 mil''). As convenções informadas na descrição da ferramenta valem como regra (``detalhada'' $\to$ 1:25.000; ``pequena escala'' $\to$ 1:250.000). Expressões qualitativas não informadas ao modelo aceitam o conjunto de escalas que as satisfazem: ``grande escala'' $\to$ 1:25.000 ou maior; ``média escala'' $\to$ 1:50.000 ou 1:100.000; ``maior que 100k'' $\to$ 1:50.000 ou maior. Duas escalas explícitas: aceita-se qualquer uma.
  \item[\texttt{productType}] Sinônimos informados na ferramenta (``topo'' e ``topográfica''; ``orto'' e ``ortoimg''; ``mdt'') e demais tipos pelo nome. ``Cartas'', ``mapas'', ``produtos'' e ``folhas'' sozinhos não definem tipo.
  \item[\texttt{state}, \texttt{city}] Siglas e grafias sem acento são normalizadas. ``Rio de Janeiro'', ``São Paulo'' e ``rio'' sem qualificador admitem as leituras estado ou município; siglas são sempre estado. Regiões (``sul'', ``Nordeste'', ``Amazônia Legal'') não são estados e não são anotadas, pois o \textit{schema} não tem campo de região.
  \item[\texttt{supplyArea}] Qualquer menção ao Centro de Geoinformação (``1º CGEO'', ``1o cgeo'', ``primeiro cgeo'', ``dois cgeo'').
  \item[\texttt{project}] Pelo nome ou pelos apelidos informados na ferramenta, com ou sem acento. A menção explícita a um projeto por um termo que o identifica sem ambiguidade no enumerado também vale como nome (``projeto rondonia'' e ``BCD de Rondônia'', sigla de Base Cartográfica Digital, para \texttt{Base Cartográfica Digital de Rondônia}). Expressão que designa o projeto só pelo termo distintivo do nome, sem a palavra ``projeto'', sem sigla do nome e sem apelido informado (``cartografia sistemática'' para \texttt{Mapeamento Sistemático}), torna o campo opcional: a consulta sustenta a leitura, mas a ferramenta não a informa.
  \item[Períodos] Verbo de publicação (``publicad-'') decide \texttt{publicationPeriod}; verbo de criação (``criad-'', ``feit-'', ``elaborad-'', ``produzid-'') decide \texttt{creationPeriod}; sem verbo que desempate, a leitura preferencial é \texttt{publicationPeriod} e \texttt{creationPeriod} é leitura estrutural alternativa, com o mesmo intervalo. O intervalo é resolvido com a data de referência da execução --- a mesma informada ao modelo --- segundo o Quadro~\ref{quadro:tempo-relativo}.
  \item[Ordenação] ``Mais recente(s)'', ``último(s)'' $\to$ \texttt{DESC}; ``mais antiga(s)'', ``primeiro(s)'', ``ordem cronológica'' $\to$ \texttt{ASC}. Pista de criação ligada à ordenação (``primeiro mapeamento feito'', ``última atualização'', ``criadas mais recentemente'') $\to$ \texttt{creationDate}; pista de publicação (``últimas publicações'', ``última carta publicada'') $\to$ \texttt{publicationDate}; sem pista, aceita-se qualquer dos dois campos. Número explícito $\to$ \texttt{limit}; singular sem número (``a carta mais recente'') $\to$ \texttt{limit = 1} opcional; plural sem número $\to$ sem \texttt{limit}. Ordenação nunca é anotada por padrão.
\end{description}

\begin{quadro}[htbp!]
\centering
\caption{Leituras aceitas para as expressões de tempo relativo (D = data de referência da execução)}
\label{quadro:tempo-relativo}
\small
\begin{tabular}{|p{4.6cm}|p{9.4cm}|}
\hline
\textbf{Expressão} & \textbf{Leituras aceitas (a primeira é a preferencial)} \\
\hline
esse ano / este ano / deste ano & ano inteiro; 1º de janeiro até D \\ \hline
ano passado & ano anterior inteiro \\ \hline
2 anos atrás & o ano de dois anos antes, inteiro; D menos 2 anos até D \\ \hline
mês passado & mês anterior inteiro \\ \hline
este mês / deste mês & dia 1 até D; mês inteiro \\ \hline
semana passada & D$-$7 até D$-$1; D$-$7 até D; semana-calendário anterior (segunda a domingo) \\ \hline
esta semana / nesta semana & segunda-feira até D; segunda a domingo \\ \hline
hoje & D a D \\ \hline
últimos 3 (6) meses & D$-$90 (180) dias até D; mesmo dia, 3 (6) meses antes, até D \\ \hline
últimos 5 anos & D menos 5 anos até D; 1º de janeiro do ano$-$5 até D \\ \hline
desde AAAA & AAAA-01-01 até D; AAAA-01-01 sem limite final \\ \hline
depois de AAAA & início em AAAA-01-01 ou (AAAA+1)-01-01; fim em D ou sem limite final \\ \hline
antes de AAAA & sem limite inicial até (AAAA$-$1)-12-31 \\ \hline
primeiro trimestre (deste ano) & janeiro a março do ano de D \\ \hline
segundo semestre do ano passado & julho a dezembro do ano anterior \\ \hline
último trimestre do ano passado & outubro a dezembro do ano anterior \\ \hline
trimestre passado & trimestre-calendário anterior ao de D \\ \hline
\end{tabular}
\fonte{Elaborado pelos autores. Implementação: \texttt{relative\_time.py} do repositório.}
\end{quadro}

\subsection*{Consultas fora do domínio}

A ferramenta não deve ser chamada quando a consulta trata de outro assunto (clima, preço, ficção), pede território fora do Brasil ou fora da Terra, ou não tem intenção de busca. Essas consultas formam a categoria F (fora do domínio) e entram nas métricas principais: a resposta é correta quando a ferramenta não é chamada. Consultas com valor impossível de representar (``escala 1:10.000.000'', ``estado 42'') ou cujo comportamento esperado é discutível (``cartas do futuro'', ``cartas topo no oceano atlântico'') são observacionais.

\subsection*{Rastreabilidade}

Cada caso registra a origem. Na camada P, a fonte é o arquivo de testes do protótipo (\texttt{backend/evaluation/test-cases.ts}), com a linha da consulta; a anotação original da equipe do 1º CGEO é mantida como leitura preferencial, e as leituras alternativas acrescentadas seguem este manual. Na camada N, a consulta foi redigida pelos autores, e a justificativa de cada leitura alternativa fica registrada em nota. Na camada G, a consulta foi gerada pelo \textit{script} gerador (semente 42), que registra a função geradora e o modelo de frase; o gabarito é conhecido por construção.

\subsection*{Procedimento de auditoria}

\begin{enumerate}
  \item \textbf{Checagens automáticas} (\texttt{pfc-auditar}): consultas duplicadas ou contraditórias; validade de enumerados e campos; resolução de todas as expressões de tempo em datas diversas; e um anotador baseado em regras, independente do gabarito, que localiza na consulta o trecho que justifica cada campo e aponta campo anotado sem evidência, evidência sem campo anotado e valor divergente.
  \item \textbf{Anotação independente às cegas}: cada consulta é anotada novamente, sem acesso ao gabarito, a partir apenas deste manual e da definição da ferramenta. A concordância é medida em rigor crescente: decisão (chamar, não chamar ou observacional, com kappa de Cohen); leituras compatíveis nos dois sentidos; leitura preferencial idêntica; presença e valor de cada campo; e detecção de ambiguidade, também com kappa de Cohen.
  \item \textbf{Adjudicação}: cada consulta com divergência em qualquer das medidas recebe uma decisão registrada --- gabarito mantido, corrigido ou ampliado com leitura alternativa ---, com justificativa por referência a este manual e com o gabarito anterior à decisão, a partir do qual a concordância anterior à adjudicação é recalculada.
  \item O processo se repete até não restar divergência sem decisão.
  \item \textbf{Revisão humana de amostra}: uma amostra estratificada por origem (40 consultas, semente 42) é conferida por uma pessoa, que marca se concorda com as leituras aceitas.
\end{enumerate}


\section{Consultas do protótipo (P01--P22)}
\label{apx:dataset-p}

As 22 consultas do arquivo de testes do protótipo do 1º CGEO (\texttt{backend/\allowbreak evaluation/\allowbreak test-cases.ts}) aparecem na ordem do arquivo. A anotação original da equipe do 1º CGEO é a leitura preferencial, e as leituras alternativas acrescentadas seguem o manual. No arquivo original, os períodos relativos estavam fixados em datas calculadas para o início de 2025 (por exemplo, 2024-10-01 a 2024-12-31 para o ``último trimestre do ano passado'' de P15). Aqui, eles são expressos pelas regras nomeadas do Quadro~\ref{quadro:tempo-relativo}.

% Nas listagens, a aspa reta é literal: sem isto, o atalho " do babel (brazil) engole o espaço seguinte.
\shorthandoff{"}
% AUTO-GERADO por pfc-dataset --apendice a partir de manual_cases.py — não editar à mão

\dtcasenota{P01}{M · C}{preciso da carta MI 2965-2-NE do segundo cgeo}{keyword: "2965-2-NE"\newline supplyArea: "2$^\circ$ Centro de Geoinformação"}{origem: test-cases.ts do protótipo, linha 6}
\dtcasenota{P02}{M · C}{folha MI 2866-3 em escala 50k}{keyword: "2866-3"\newline scale: "1:50.000"}{origem: test-cases.ts do protótipo, linha 16}
\dtcasenota{P03}{M · C · A}{MI 2901 do projeto olimpiadas}{keyword: "2901"\newline project: "Olimpíadas Rio 2016"}{origem: test-cases.ts do protótipo, linha 26}
\dtcasenota{P04}{M · C · A}{carta 530 em pequena escala}{keyword: "530"\newline scale: "1:250.000"}{origem: test-cases.ts do protótipo, linha 36}
\dtcasenota{P05}{M · C · T}{buscar folha SF-22-Y-D-II-4 do RS publicada esse ano}{keyword: "SF-22-Y-D-II-4"\newline state: "Rio Grande do Sul"\newline publicationPeriod: ano corrente}{origem: test-cases.ts do protótipo, linha 48}
\dtcasenota{P06}{M · C · A}{carta SF-22-Y-D do tipo ortoimg}{keyword: "SF-22-Y-D"\newline productType: "SCN Carta Ortoimagem"}{origem: test-cases.ts do protótipo, linha 62}
\dtcasenota{P07}{M · C · A}{inom SF-22-Y-D-II-4-SE criado pelo primeiro cgeo}{keyword: "SF-22-Y-D-II-4-SE"\newline supplyArea: "1$^\circ$ Centro de Geoinformação"}{origem: test-cases.ts do protótipo, linha 72}
\dtcasenota{P08}{C}{carta topográfica Passo da Seringueira em 1:25.000}{keyword: "Passo da Seringueira"\newline scale: "1:25.000"\newline productType: "SCN Carta Topográfica Matricial"}{origem: test-cases.ts do protótipo, linha 84}
\dtcasenota{P09}{C}{folha Porto Velho publicada em 2024}{keyword: "Porto Velho"\newline publicationPeriod: 2024-01-01 a 2024-12-31\newline aceita-se também city no lugar de keyword}{origem: test-cases.ts do protótipo, linha 95; nome de carta coincide com município: aceita-se também city}
\dtcasenota{P10}{C · A}{carta Vale do Guaporé do projeto rondonia}{keyword: "Vale do Guaporé"\newline project: "Base Cartográfica Digital de Rondônia"}{origem: test-cases.ts do protótipo, linha 108}
\dtcasenota{P11}{M · C · A}{MI 2965-2-NE e SF-22-Y-D do terceiro cgeo em grande escala}{keyword: "2965-2-NE" ou "SF-22-Y-D"\newline supplyArea: "3$^\circ$ Centro de Geoinformação"\newline scale: "1:25.000" ou "1:10.000" ou "1:5.000" ou "1:2.000" ou "1:1.000"}{origem: test-cases.ts do protótipo, linha 120; dois códigos (o schema comporta um): aceita-se qualquer um; 'grande escala' aceita as escalas grandes}
\dtcasenota{P12}{M · C · T · A}{cartas do tipo topo com MI 530 ou SF-22 publicadas esse ano}{keyword: "530" ou "SF-22"\newline productType: "SCN Carta Topográfica Matricial"\newline publicationPeriod: ano corrente}{origem: test-cases.ts do protótipo, linha 131; dois códigos: aceita-se qualquer um}
\dtcasenota{P13}{C · A}{todas as cartas em escala maior que 100k do rio grande do sul}{state: "Rio Grande do Sul"\newline scale: "1:25.000" ou "1:50.000" ou "1:10.000" ou "1:5.000" ou "1:2.000" ou "1:1.000"}{origem: test-cases.ts do protótipo, linha 147; 'maior que 100k' aceita qualquer escala maior que 1:100.000}
\dtcasenota{P14}{C · A}{mapeamento detalhado (25k ou 50k) de brasilia}{scale: "1:25.000" ou "1:50.000"\newline city: "Brasília"}{origem: test-cases.ts do protótipo, linha 156; duas escalas explícitas: aceita-se qualquer uma}
\dtcasenota{P15}{T · C · A}{cartas criadas no último trimestre do ano passado do segundo cgeo}{supplyArea: "2$^\circ$ Centro de Geoinformação"\newline creationPeriod: 4º trimestre do ano anterior}{origem: test-cases.ts do protótipo, linha 167}
\dtcasenota{P16}{T · C · A}{produtos da semana passada do terceiro cgeo}{supplyArea: "3$^\circ$ Centro de Geoinformação"\newline publicationPeriod: semana passada\newline aceita-se também creationPeriod no lugar de publicationPeriod}{origem: test-cases.ts do protótipo, linha 179; sem verbo que decida publicação ou criação}
\dtcasenota{P17}{C · A}{cartas do mapeamento sistematico em goias}{project: "Mapeamento Sistemático"\newline state: "Goiás"}{origem: test-cases.ts do protótipo, linha 193}
\dtcasenota{P18}{C · A}{produtos da copa das confederacoes em fortaleza}{project: "Copa das Confederações"\newline city: "Fortaleza"}{origem: test-cases.ts do protótipo, linha 202}
\dtcasenota{P19}{O · C · A}{primeiro mapeamento feito em cuiaba}{city: "Cuiabá"\newline sortField: "creationDate"\newline sortDirection: "ASC"\newline limit: 1 (opcional)}{origem: test-cases.ts do protótipo, linha 213; singular sem número: limit 1 opcional}
\dtcasenota{P20}{O · C · A}{ultima atualizacao do para}{state: "Pará"\newline sortField: "creationDate"\newline sortDirection: "DESC"\newline limit: 1 (opcional)}{origem: test-cases.ts do protótipo, linha 223; singular sem número: limit 1 opcional}
\dtcasenota{P21}{O · T · C · A}{preciso das 5 cartas mais antigas do tipo ortoimg em escala detalhada do terceiro cgeo em pernambuco publicadas depois de 2020}{limit: 5\newline productType: "SCN Carta Ortoimagem"\newline scale: "1:25.000"\newline supplyArea: "3$^\circ$ Centro de Geoinformação"\newline state: "Pernambuco"\newline publicationPeriod: depois de 2020\newline sortField: "creationDate" ou "publicationDate"\newline sortDirection: "ASC"}{origem: test-cases.ts do protótipo, linha 235; 'mais antigas' sem pista ligada à ordenação: aceita-se publicationDate além do creationDate original}
\dtcasenota{P22}{M · O · T · C}{MI 2901 ou SF-22 do quarto cgeo no sudeste em pequena escala criado entre 2022 e 2023 ordem cronologica}{keyword: "2901" ou "SF-22"\newline supplyArea: "4$^\circ$ Centro de Geoinformação"\newline scale: "1:250.000"\newline creationPeriod: 2022-01-01 a 2023-12-31\newline sortField: "creationDate" ou "publicationDate"\newline sortDirection: "ASC"}{origem: test-cases.ts do protótipo, linha 253; dois códigos; 'sudeste' é região (nada a anotar); ordem cronológica sem pista ligada}


\section{Consultas redigidas pelos autores (N01--N40)}
\label{apx:dataset-n}

Os autores redigiram estas consultas para preencher lacunas de distribuição, em especial a categoria Simples, ausente na camada P. Elas exercitam também as ambiguidades previstas no manual (estado ou capital, ordenação sem pista, regiões e grafias informais de escala, como ``25k'' e ``cem k'') e a fronteira do domínio. N36 e N37 estão fora do domínio (categoria F), e N38 a N40 são observacionais.

% AUTO-GERADO por pfc-dataset --apendice a partir de manual_cases.py — não editar à mão

\dtcasenota{N01}{S}{cartas de São Paulo}{state: "São Paulo"\newline aceita-se também city no lugar de state}{'São Paulo' sem qualificador: estado ou capital}
\dtcase{N02}{S}{produtos em escala 1:50.000}{scale: "1:50.000"}
\dtcase{N03}{S}{ortoimagens}{productType: "SCN Carta Ortoimagem"}
\dtcase{N04}{S · A}{mapas do rj}{state: "Rio de Janeiro"}
\dtcase{N05}{S · A}{carta 25k}{scale: "1:25.000"}
\dtcase{N06}{S}{folha SF-22-Y-D}{keyword: "SF-22-Y-D"}
\dtcase{N07}{S · A}{modelos digitais de terreno}{productType: "MDT — RAM"}
\dtcase{N08}{S · A}{cartas topográficas}{productType: "SCN Carta Topográfica Matricial"}
\dtcase{N09}{S}{produtos do Amazonas}{state: "Amazonas"}
\dtcase{N10}{S · A}{mi 2901}{keyword: "2901"}
\dtcase{N11}{C}{cartas do 1º CGEO em Minas Gerais}{supplyArea: "1$^\circ$ Centro de Geoinformação"\newline state: "Minas Gerais"}
\dtcase{N12}{C}{ortoimagens de Manaus em 1:50.000}{productType: "SCN Carta Ortoimagem"\newline city: "Manaus"\newline scale: "1:50.000"}
\dtcase{N13}{C · A}{mdt do para em 25k}{productType: "MDT — RAM"\newline state: "Pará"\newline scale: "1:25.000"}
\dtcase{N14}{M · C}{folha 2965-2 do Rio Grande do Sul}{keyword: "2965-2"\newline state: "Rio Grande do Sul"}
\dtcase{N15}{M · C}{INOM SG-22-X-A}{keyword: "SG-22-X-A"}
\dtcase{N16}{M · C · T}{carta MI 3010 publicada em 2023}{keyword: "3010"\newline publicationPeriod: 2023-01-01 a 2023-12-31}
\dtcase{N17}{T · A}{produtos publicados nos últimos 3 meses}{publicationPeriod: últimos 3 meses}
\dtcasenota{N18}{T · C}{cartas do ano passado em Roraima}{state: "Roraima"\newline publicationPeriod: ano anterior, inteiro\newline aceita-se também creationPeriod no lugar de publicationPeriod}{sem verbo que decida publicação ou criação}
\dtcasenota{N19}{T · A}{atualizações deste mês}{publicationPeriod: mês corrente\newline aceita-se também creationPeriod no lugar de publicationPeriod}{sem verbo que decida publicação ou criação}
\dtcasenota{N20}{T · C · A}{cartas dos últimos 5 anos em escala 100k}{scale: "1:100.000"\newline publicationPeriod: últimos 5 anos\newline aceita-se também creationPeriod no lugar de publicationPeriod}{sem verbo que decida publicação ou criação}
\dtcasenota{N21}{T · A}{qualquer coisa de hoje}{publicationPeriod: dia da execução\newline aceita-se também creationPeriod no lugar de publicationPeriod}{sem verbo que decida publicação ou criação}
\dtcasenota{N22}{T · C · A}{ortoimagens de 2 anos atrás}{productType: "SCN Carta Ortoimagem"\newline publicationPeriod: ano de dois anos antes\newline aceita-se também creationPeriod no lugar de publicationPeriod}{sem verbo que decida publicação ou criação}
\dtcasenota{N23}{T · C · A}{produtos do trimestre passado do 4º CGEO}{supplyArea: "4$^\circ$ Centro de Geoinformação"\newline publicationPeriod: trimestre anterior\newline aceita-se também creationPeriod no lugar de publicationPeriod}{sem verbo que decida publicação ou criação}
\dtcasenota{N24}{O · C}{3 cartas mais recentes de Curitiba}{city: "Curitiba"\newline limit: 3\newline sortField: "publicationDate" ou "creationDate"\newline sortDirection: "DESC"}{'mais recentes' sem pista: aceita-se publicationDate ou creationDate}
\dtcasenota{N25}{O · C}{10 primeiras ortoimagens do Nordeste}{productType: "SCN Carta Ortoimagem"\newline limit: 10\newline sortField: "publicationDate" ou "creationDate"\newline sortDirection: "ASC"}{'Nordeste' é região (nada a anotar); 'primeiras' sem pista}
\dtcasenota{N26}{O · C · A}{mais antigas em rondonia}{state: "Rondônia"\newline sortField: "publicationDate" ou "creationDate"\newline sortDirection: "ASC"}{'mais antigas' sem pista}
\dtcase{N27}{O · C}{últimas 5 publicações}{limit: 5\newline sortField: "publicationDate"\newline sortDirection: "DESC"}
\dtcasenota{N28}{O · M · C}{carta MI 2901 mais recente}{keyword: "2901"\newline sortField: "publicationDate" ou "creationDate"\newline sortDirection: "DESC"\newline limit: 1 (opcional)}{singular sem número: limit 1 opcional; 'mais recente' sem pista}
\dtcasenota{N29}{O · C · T}{três produtos mais antigos deste ano}{limit: 3\newline publicationPeriod: ano corrente\newline sortField: "publicationDate" ou "creationDate"\newline sortDirection: "ASC"\newline aceita-se também creationPeriod no lugar de publicationPeriod}{sem verbo que decida o período; 'mais antigos' sem pista}
\dtcase{N30}{A · C}{cartas topo do primeiro cgeo}{productType: "SCN Carta Topográfica Matricial"\newline supplyArea: "1$^\circ$ Centro de Geoinformação"}
\dtcase{N31}{A · S}{orto banda p}{productType: "SCN Carta Ortoimagem Banda P Pol HH"}
\dtcasenota{N32}{A · C}{mapas do rio em cem k}{state: "Rio de Janeiro"\newline scale: "1:100.000"\newline aceita-se também city no lugar de state}{'rio' é ambíguo entre o estado e a capital}
\dtcasenota{N33}{A · C}{cartografia sistematica sp}{project: "Mapeamento Sistemático" (opcional)\newline state: "São Paulo"}{'cartografia sistemática' designa o programa pelo termo distintivo do nome, sem ser nome nem apelido informado na ferramenta: project é opcional (regra de project; adjudicação)}
\dtcasenota{N34}{A · C}{ortos da amazonia legal}{productType: "SCN Carta Ortoimagem"}{'Amazônia Legal' é região (nove estados), não estado: nada a anotar para ela}
\dtcase{N35}{A · M · C}{sf22yd do dois cgeo}{keyword: "SF-22-Y-D"\newline supplyArea: "2$^\circ$ Centro de Geoinformação"}
\dtcase{N36}{F}{cartas fora do brasil}{não chamar a ferramenta (fora do território brasileiro, P4)}
\dtcase{N37}{F}{quero um mapa bonito}{não chamar a ferramenta (sem critério de busca, P4)}
\dtcase{N38}{S · A}{escala 1:10.000.000}{não chamar a ferramenta (observacional, P6: escala fora do enumerado da ferramenta)}
\dtcasenota{N39}{C · A}{cartas topo no oceano atlântico}{productType: "SCN Carta Topográfica Matricial"\newline caso observacional (P6): executado e reportado, fora das métricas principais}{área não terrestre: comportamento esperado discutível}
\dtcasenota{N40}{C · T}{cartas do futuro (ano 2050)}{publicationPeriod: 2050-01-01 a 2050-12-31\newline aceita-se também creationPeriod no lugar de publicationPeriod\newline caso observacional (P6): executado e reportado, fora das métricas principais}{data futura: comportamento esperado discutível}


\section{Geração programática da camada G}
\label{sec:dataset-gerado}

O \textit{script} \texttt{scripts/\allowbreak generate\_dataset\_paper.py} do repositório (Apêndice~\ref{apendice:codigo}) produz as 248 consultas da camada G a partir de catálogos fechados do domínio. São as 27 unidades da federação (com preposição, grafias usuais e, exceto para Sergipe, sigla), 12 municípios de referência sem homônimo entre as UFs, as 8 escalas do SCN (com apelidos como ``25k'', ``25 mil'' e as convenções ``detalhada'' e ``pequena escala''), os 9 tipos de produto e os 9 projetos institucionais aceitos pela ferramenta, os 5 CGEOs, 12 códigos MI e 10 códigos INOM, 15 expressões de tempo (14 relativas e o intervalo absoluto ``entre 2022 e 2023'') e os padrões de ordenação. Oito funções geradoras, uma por família, combinam esses catálogos em modelos de frase do português corrente.

Três propriedades tornam a camada auditável. A geração é \textbf{determinística}, porque um único gerador pseudoaleatório, com semente 42, é consumido em ordem fixa, e a mesma versão do \textit{script} produz sempre as mesmas consultas. O gabarito é conhecido \textbf{por construção}, porque a frase é montada a partir dos parâmetros que deve produzir, e cada consulta registra no \textit{dataset} a função geradora e o modelo de frase que a originaram. \textbf{Nenhuma consulta se repete} depois da normalização de caixa, acentos e pontuação, dentro da camada e em relação às camadas P e N. Valem as convenções do manual, como as leituras estado ou município para ``São Paulo'' e ``Rio de Janeiro'' por extenso e os dois campos de período nas frases sem verbo de publicação ou criação.

As seções seguintes listam as 248 consultas por família. Os identificadores seguem o padrão \texttt{G\{família\}\{sequência\}}, em que a família é S (simples), C (compostas), M (código MI/INOM), T (tempo relativo), O (ordenação), A (informais), P (amplas, com muitos produtos) ou F (fronteira do domínio).

% ─── AUTO-GERADO por generate_dataset.py (semente 42): 248 consultas ───

\section{Consultas simples geradas por \textit{templates} (G-S)}

Total desta fam\'ilia: \textbf{61 consultas}, listadas integralmente abaixo (a fun\c{c}\~ao geradora e o modelo de frase de cada consulta constam de \texttt{dataset.json}).

{\footnotesize
\begin{longtable}{|p{1.1cm}|p{1.3cm}|p{5.4cm}|p{6.3cm}|}
\hline
\textbf{ID} & \textbf{Cat.} & \textbf{Consulta} & \textbf{Leituras aceitas} \\ \hline
\endhead
GS001 & S & cartas do Acre & state: "Acre" \\ \hline
GS002 & S · A & cartas de AL & state: "Alagoas" \\ \hline
GS003 & S · A & cartas do AP & state: "Amapá" \\ \hline
GS004 & S & cartas do Amazonas & state: "Amazonas" \\ \hline
GS005 & S · A & cartas da bahia & state: "Bahia" \\ \hline
GS006 & S · A & cartas do CE & state: "Ceará" \\ \hline
GS007 & S · A & cartas do DF & state: "Distrito Federal" \\ \hline
GS008 & S · A & cartas do ES & state: "Espírito Santo" \\ \hline
GS009 & S & cartas de Goiás & state: "Goiás" \\ \hline
GS010 & S · A & cartas do MA & state: "Maranhão" \\ \hline
GS011 & S & cartas de Mato Grosso & state: "Mato Grosso" \\ \hline
GS012 & S & cartas de Mato Grosso do Sul & state: "Mato Grosso do Sul" \\ \hline
GS013 & S · A & cartas de MG & state: "Minas Gerais" \\ \hline
GS014 & S & cartas do Pará & state: "Pará" \\ \hline
GS015 & S · A & cartas da paraíba & state: "Paraíba" \\ \hline
GS016 & S · A & produtos na escala 1:1000 & scale: "1:1.000" \\ \hline
GS017 & S · A & produtos na escala 1:2000 & scale: "1:2.000" \\ \hline
GS018 & S · A & produtos na escala 1:5000 & scale: "1:5.000" \\ \hline
GS019 & S · A & produtos na escala 1:10000 & scale: "1:10.000" \\ \hline
GS020 & S · A & produtos na escala 1:25000 & scale: "1:25.000" \\ \hline
GS021 & S · A & produtos na escala de 50 mil & scale: "1:50.000" \\ \hline
GS022 & S · A & produtos na escala de 100 mil & scale: "1:100.000" \\ \hline
GS023 & S · A & produtos em 250k & scale: "1:250.000" \\ \hline
GS024 & S · A & cartas topográficas disponíveis & productType: "SCN Carta Topográfica Matricial" \\ \hline
GS025 & S · A & cartas vetoriais disponíveis & productType: "SCN Carta Topográfica Vetorial" \\ \hline
GS026 & S · A & ortoimagens disponíveis & productType: "SCN Carta Ortoimagem" \\ \hline
GS027 & S · A & ortoimagens banda P HH disponíveis & productType: "SCN Carta Ortoimagem Banda P Pol HH" \\ \hline
GS028 & S · A & ortoimagens banda X HH disponíveis & productType: "SCN Carta Ortoimagem Banda X Pol HH" \\ \hline
GS029 & S · A & MDTs disponíveis & productType: "MDT — RAM" \\ \hline
GS030 & S · A & MDSs disponíveis & productType: "MDS — RAM" \\ \hline
GS031 & S · A & cartas CIRC disponíveis & productType: "CIRC" \\ \hline
GS032 & S · A & mapas temáticos disponíveis & productType: "Cartas Temáticas Não SCN" \\ \hline
GS033 & S · A & produtos do primeiro cgeo & supplyArea: "1$^\circ$ Centro de Geoinformação" \\ \hline
GS034 & S · A & produtos do 2o cgeo & supplyArea: "2$^\circ$ Centro de Geoinformação" \\ \hline
GS035 & S · A & produtos do 3o cgeo & supplyArea: "3$^\circ$ Centro de Geoinformação" \\ \hline
GS036 & S · A & produtos do quarto cgeo & supplyArea: "4$^\circ$ Centro de Geoinformação" \\ \hline
GS037 & S · A & produtos do 5º CGEO & supplyArea: "5$^\circ$ Centro de Geoinformação" \\ \hline
GS038 & S & cartas de Brasília & city: "Brasília" \\ \hline
GS039 & S & cartas de Manaus & city: "Manaus" \\ \hline
GS040 & S & cartas de Porto Velho & city: "Porto Velho" \\ \hline
GS041 & S & cartas de Cuiabá & city: "Cuiabá" \\ \hline
GS042 & S & cartas de Fortaleza & city: "Fortaleza" \\ \hline
GS043 & S & cartas de Salvador & city: "Salvador" \\ \hline
GS044 & S & cartas de Belo Horizonte & city: "Belo Horizonte" \\ \hline
GS045 & S & cartas de Curitiba & city: "Curitiba" \\ \hline
GS046 & S & cartas de Porto Alegre & city: "Porto Alegre" \\ \hline
GS047 & S & cartas de Recife & city: "Recife" \\ \hline
GS048 & S · M & carta MI 2965-2-NE & keyword: "2965-2-NE" \\ \hline
GS049 & S · M & carta MI 2866-3 & keyword: "2866-3" \\ \hline
GS050 & S · M & carta MI 2901 & keyword: "2901" \\ \hline
GS051 & S · M & carta MI 530 & keyword: "530" \\ \hline
GS052 & S · M & carta MI 3010 & keyword: "3010" \\ \hline
GS053 & S · M & carta MI 2965 & keyword: "2965" \\ \hline
GS054 & S · M & carta MI 2866-2-SO & keyword: "2866-2-SO" \\ \hline
GS055 & S · M & carta MI 2901-4-NE & keyword: "2901-4-NE" \\ \hline
GS056 & S · M & folha INOM SF-22-Y-D & keyword: "SF-22-Y-D" \\ \hline
GS057 & S · M & folha INOM SF-22-Y-D-II & keyword: "SF-22-Y-D-II" \\ \hline
GS058 & S · M & folha INOM SF-22-Y-D-II-4 & keyword: "SF-22-Y-D-II-4" \\ \hline
GS059 & S · M & folha INOM SF-22-Y-D-II-4-SE & keyword: "SF-22-Y-D-II-4-SE" \\ \hline
GS060 & S · M & folha INOM SG-22-X-A & keyword: "SG-22-X-A" \\ \hline
GS061 & S · M & folha INOM SF-23-V-C & keyword: "SF-23-V-C" \\ \hline
\end{longtable}
\addtocounter{table}{-1}
}

\section{Consultas compostas geradas por \textit{templates} (G-C)}

Total desta fam\'ilia: \textbf{55 consultas}, listadas integralmente abaixo (a fun\c{c}\~ao geradora e o modelo de frase de cada consulta constam de \texttt{dataset.json}).

{\footnotesize
\begin{longtable}{|p{1.1cm}|p{1.3cm}|p{5.4cm}|p{6.3cm}|}
\hline
\textbf{ID} & \textbf{Cat.} & \textbf{Consulta} & \textbf{Leituras aceitas} \\ \hline
\endhead
GC001 & C · A & cartas do amapá na escala 1:2.000 & state: "Amapá"; scale: "1:2.000" \\ \hline
GC002 & C · A & cartas de Mato Grosso do Sul em 25k & state: "Mato Grosso do Sul"; scale: "1:25.000" \\ \hline
GC003 & C · A & cartas de santa catarina na escala 1:2.000 & state: "Santa Catarina"; scale: "1:2.000" \\ \hline
GC004 & C · A & cartas do Amapá na escala de 25 mil & state: "Amapá"; scale: "1:25.000" \\ \hline
GC005 & C · A & cartas do RJ na escala 1:2000 & state: "Rio de Janeiro"; scale: "1:2.000" \\ \hline
GC006 & C · A & cartas de RO em 25k & state: "Rondônia"; scale: "1:25.000" \\ \hline
GC007 & C · A & cartas do ES na escala 1:100000 & state: "Espírito Santo"; scale: "1:100.000" \\ \hline
GC008 & C · A & cartas da Paraíba em 50k & state: "Paraíba"; scale: "1:50.000" \\ \hline
GC009 & C · A & cartas de mato grosso do sul na escala de 10 mil & state: "Mato Grosso do Sul"; scale: "1:10.000" \\ \hline
GC010 & C · A & cartas de Goiás na escala 1:2000 & state: "Goiás"; scale: "1:2.000" \\ \hline
GC011 & C · A & cartas do Piauí na escala 1:10000 & state: "Piauí"; scale: "1:10.000" \\ \hline
GC012 & C · A & cartas da paraíba na escala 1:25000 & state: "Paraíba"; scale: "1:25.000" \\ \hline
GC013 & C · A & cartas de rondônia na escala 1:1000 & state: "Rondônia"; scale: "1:1.000" \\ \hline
GC014 & C · A & cartas do TO na escala 1:50000 & state: "Tocantins"; scale: "1:50.000" \\ \hline
GC015 & C · A & cartas de GO na escala de 10 mil & state: "Goiás"; scale: "1:10.000" \\ \hline
GC016 & C · A & cartas de roraima na escala de 10 mil & state: "Roraima"; scale: "1:10.000" \\ \hline
GC017 & C · A & cartas do paraná na escala 1:250000 & state: "Paraná"; scale: "1:250.000" \\ \hline
GC018 & C · A & cartas do Piauí em escala detalhada & state: "Piauí"; scale: "1:25.000" \\ \hline
GC019 & C · A & cartas do rio de janeiro em 50k & state: "Rio de Janeiro"; scale: "1:50.000"; aceita-se também city no lugar de state \\ \hline
GC020 & C · A & cartas da Bahia em 250k & state: "Bahia"; scale: "1:250.000" \\ \hline
GC021 & C · A & cartas de SP na escala 1:2000 & state: "São Paulo"; scale: "1:2.000" \\ \hline
GC022 & C · A & cartas do RS na escala de 100 mil & state: "Rio Grande do Sul"; scale: "1:100.000" \\ \hline
GC023 & C · A & cartas do amapá na escala de 100 mil & state: "Amapá"; scale: "1:100.000" \\ \hline
GC024 & C · A & cartas da Paraíba em 25k & state: "Paraíba"; scale: "1:25.000" \\ \hline
GC025 & C & cartas de Rondônia na escala 1:2.000 & state: "Rondônia"; scale: "1:2.000" \\ \hline
GC026 & C · A & cartas de São Paulo em 50k & state: "São Paulo"; scale: "1:50.000"; aceita-se também city no lugar de state \\ \hline
GC027 & C · A & cartas do maranhão na escala 1:5.000 & state: "Maranhão"; scale: "1:5.000" \\ \hline
GC028 & C · A & cartas do Acre na escala 1:25000 & state: "Acre"; scale: "1:25.000" \\ \hline
GC029 & C · A & cartas de PE na escala 1:25000 & state: "Pernambuco"; scale: "1:25.000" \\ \hline
GC030 & C · A & cartas da bahia na escala 1:5000 & state: "Bahia"; scale: "1:5.000" \\ \hline
GC031 & C · A & mapas de Recife na escala 1:50000 & city: "Recife"; scale: "1:50.000" \\ \hline
GC032 & C & mapas de Brasília na escala 1:2.000 & city: "Brasília"; scale: "1:2.000" \\ \hline
GC033 & C · A & mapas de Fortaleza na escala 1:10000 & city: "Fortaleza"; scale: "1:10.000" \\ \hline
GC034 & C · A & mapas de Cuiabá na escala 1:2000 & city: "Cuiabá"; scale: "1:2.000" \\ \hline
GC035 & C · A & mapas de Belém em 250k & city: "Belém"; scale: "1:250.000" \\ \hline
GC036 & C · A & mapas de Porto Alegre na escala 1:5000 & city: "Porto Alegre"; scale: "1:5.000" \\ \hline
GC037 & C · A & mapas de Goiânia na escala 1:250000 & city: "Goiânia"; scale: "1:250.000" \\ \hline
GC038 & C · A & mapas de Fortaleza em 100k & city: "Fortaleza"; scale: "1:100.000" \\ \hline
GC039 & C · A & cartas temáticas de São Paulo & productType: "Cartas Temáticas Não SCN"; state: "São Paulo"; aceita-se também city no lugar de state \\ \hline
GC040 & C · A & ortoimagens banda X HH de minas & productType: "SCN Carta Ortoimagem Banda X Pol HH"; state: "Minas Gerais" \\ \hline
GC041 & C · A & cartas temáticas da PB & productType: "Cartas Temáticas Não SCN"; state: "Paraíba" \\ \hline
GC042 & C · A & ortoimagens banda P HH do amapá & productType: "SCN Carta Ortoimagem Banda P Pol HH"; state: "Amapá" \\ \hline
GC043 & C · A & cartas temáticas do Espírito Santo & productType: "Cartas Temáticas Não SCN"; state: "Espírito Santo" \\ \hline
GC044 & C · A & cartas topográficas do Amapá & productType: "SCN Carta Topográfica Matricial"; state: "Amapá" \\ \hline
GC045 & C · A & produtos do projeto copa das confederacoes no AP & project: "Copa das Confederações"; state: "Amapá" \\ \hline
GC046 & C · A & produtos do projeto olimpiadas em PE & project: "Olimpíadas Rio 2016"; state: "Pernambuco" \\ \hline
GC047 & C & produtos do projeto NGA-BECA no Distrito Federal & project: "NGA-BECA"; state: "Distrito Federal" \\ \hline
GC048 & C · A & produtos do projeto beca no espírito santo & project: "NGA-BECA"; state: "Espírito Santo" \\ \hline
GC049 & C · A & produtos do projeto copa das confederacoes no AM & project: "Copa das Confederações"; state: "Amazonas" \\ \hline
GC050 & C · A & produtos do projeto BCD de Rondônia no pará & project: "Base Cartográfica Digital de Rondônia"; state: "Pará" \\ \hline
GC051 & C · A & cartas do 1º CGEO na escala 1:2.000 & supplyArea: "1$^\circ$ Centro de Geoinformação"; scale: "1:2.000" \\ \hline
GC052 & C · A & cartas do 3o cgeo na escala 1:2000 & supplyArea: "3$^\circ$ Centro de Geoinformação"; scale: "1:2.000" \\ \hline
GC053 & C · A & cartas do 2o cgeo em pequena escala & supplyArea: "2$^\circ$ Centro de Geoinformação"; scale: "1:250.000" \\ \hline
GC054 & C · A & cartas do 2 CGEO na escala 1:25000 & supplyArea: "2$^\circ$ Centro de Geoinformação"; scale: "1:25.000" \\ \hline
GC055 & C · A & cartas do 1º CGEO em 250k & supplyArea: "1$^\circ$ Centro de Geoinformação"; scale: "1:250.000" \\ \hline
\end{longtable}
\addtocounter{table}{-1}
}

\section{Consultas com c\'odigo MI/INOM geradas por \textit{templates} (G-M)}

Total desta fam\'ilia: \textbf{28 consultas}, listadas integralmente abaixo (a fun\c{c}\~ao geradora e o modelo de frase de cada consulta constam de \texttt{dataset.json}).

{\footnotesize
\begin{longtable}{|p{1.1cm}|p{1.3cm}|p{5.4cm}|p{6.3cm}|}
\hline
\textbf{ID} & \textbf{Cat.} & \textbf{Consulta} & \textbf{Leituras aceitas} \\ \hline
\endhead
GM001 & C · M & folha MI 2965-2-NE do Rio Grande do Sul & keyword: "2965-2-NE"; state: "Rio Grande do Sul" \\ \hline
GM002 & C · M · A & folha MI 2866-3 do TO & keyword: "2866-3"; state: "Tocantins" \\ \hline
GM003 & C · M · A & folha MI 2901 do AP & keyword: "2901"; state: "Amapá" \\ \hline
GM004 & C · M · A & folha MI 530 do ceará & keyword: "530"; state: "Ceará" \\ \hline
GM005 & C · M · A & folha MI 3010 do paraná & keyword: "3010"; state: "Paraná" \\ \hline
GM006 & C · M · A & folha MI 2965 do distrito federal & keyword: "2965"; state: "Distrito Federal" \\ \hline
GM007 & C · M · A & folha MI 2866-2-SO de AL & keyword: "2866-2-SO"; state: "Alagoas" \\ \hline
GM008 & C · M · A & folha MI 2901-4-NE de MG & keyword: "2901-4-NE"; state: "Minas Gerais" \\ \hline
GM009 & C · M · A & folha MI 531-1-NE de minas & keyword: "531-1-NE"; state: "Minas Gerais" \\ \hline
GM010 & C · M & folha MI 2965-3 de Sergipe & keyword: "2965-3"; state: "Sergipe" \\ \hline
GM011 & C · M · A & folha MI 2866-3-NO do maranhão & keyword: "2866-3-NO"; state: "Maranhão" \\ \hline
GM012 & C · M & folha MI 3010-2-SE de Roraima & keyword: "3010-2-SE"; state: "Roraima" \\ \hline
GM013 & C · M · A & INOM SF-22-Y-D do tipo temática & keyword: "SF-22-Y-D"; productType: "Cartas Temáticas Não SCN" \\ \hline
GM014 & C · M · A & INOM SF-22-Y-D-II do tipo ortoimagem & keyword: "SF-22-Y-D-II"; productType: "SCN Carta Ortoimagem" \\ \hline
GM015 & C · M · A & INOM SF-22-Y-D-II-4 do tipo banda X HH & keyword: "SF-22-Y-D-II-4"; productType: "SCN Carta Ortoimagem Banda X Pol HH" \\ \hline
GM016 & C · M · A & INOM SF-22-Y-D-II-4-SE do tipo topográfica & keyword: "SF-22-Y-D-II-4-SE"; productType: "SCN Carta Topográfica Matricial" \\ \hline
GM017 & C · M · A & INOM SG-22-X-A do tipo MDT & keyword: "SG-22-X-A"; productType: "MDT — RAM" \\ \hline
GM018 & C · M · A & INOM SF-23-V-C do tipo topo & keyword: "SF-23-V-C"; productType: "SCN Carta Topográfica Matricial" \\ \hline
GM019 & C · M · A & INOM SG-21-Z-B do tipo temática & keyword: "SG-21-Z-B"; productType: "Cartas Temáticas Não SCN" \\ \hline
GM020 & C · M · A & INOM SH-22-W-A-III-1-NE do tipo topográfica & keyword: "SH-22-W-A-III-1-NE"; productType: "SCN Carta Topográfica Matricial" \\ \hline
GM021 & C · M · A & INOM SG-22-Z-D-II-3 do tipo ortoimagem & keyword: "SG-22-Z-D-II-3"; productType: "SCN Carta Ortoimagem" \\ \hline
GM022 & C · M · A & INOM SF-24-Y-A do tipo vetorial & keyword: "SF-24-Y-A"; productType: "SCN Carta Topográfica Vetorial" \\ \hline
GM023 & C · M · A & MI 2866-2-SO produzido pelo 1 CGEO & keyword: "2866-2-SO"; supplyArea: "1$^\circ$ Centro de Geoinformação" \\ \hline
GM024 & C · M · A & MI 2866-3 produzido pelo quinto cgeo & keyword: "2866-3"; supplyArea: "5$^\circ$ Centro de Geoinformação" \\ \hline
GM025 & C · M · A & MI 2965-3 produzido pelo 3o cgeo & keyword: "2965-3"; supplyArea: "3$^\circ$ Centro de Geoinformação" \\ \hline
GM026 & C · M · A & MI 530 produzido pelo 3o cgeo & keyword: "530"; supplyArea: "3$^\circ$ Centro de Geoinformação" \\ \hline
GM027 & C · M · A & MI 2965-2-NE produzido pelo 3 CGEO & keyword: "2965-2-NE"; supplyArea: "3$^\circ$ Centro de Geoinformação" \\ \hline
GM028 & C · M · A & MI 3010 produzido pelo segundo cgeo & keyword: "3010"; supplyArea: "2$^\circ$ Centro de Geoinformação" \\ \hline
\end{longtable}
\addtocounter{table}{-1}
}

\section{Consultas com tempo relativo geradas por \textit{templates} (G-T)}

Total desta fam\'ilia: \textbf{41 consultas}, listadas integralmente abaixo (a fun\c{c}\~ao geradora e o modelo de frase de cada consulta constam de \texttt{dataset.json}).

{\footnotesize
\begin{longtable}{|p{1.1cm}|p{1.3cm}|p{5.4cm}|p{6.3cm}|}
\hline
\textbf{ID} & \textbf{Cat.} & \textbf{Consulta} & \textbf{Leituras aceitas} \\ \hline
\endhead
GT001 & T & cartas publicadas esse ano & publicationPeriod: ano corrente \\ \hline
GT002 & T & cartas publicadas no ano passado & publicationPeriod: ano anterior, inteiro \\ \hline
GT003 & T & cartas publicadas no mês passado & publicationPeriod: mês anterior, inteiro \\ \hline
GT004 & T & cartas publicadas na semana passada & publicationPeriod: semana passada \\ \hline
GT005 & T & cartas publicadas nos últimos 3 meses & publicationPeriod: últimos 3 meses \\ \hline
GT006 & T & cartas publicadas nos últimos 6 meses & publicationPeriod: últimos 6 meses \\ \hline
GT007 & T & cartas publicadas nos últimos 5 anos & publicationPeriod: últimos 5 anos \\ \hline
GT008 & T & cartas publicadas desde 2020 & publicationPeriod: desde 2020 \\ \hline
GT009 & T & cartas publicadas depois de 2022 & publicationPeriod: depois de 2022 \\ \hline
GT010 & T & cartas publicadas antes de 2020 & publicationPeriod: antes de 2020 \\ \hline
GT011 & T & cartas publicadas entre 2022 e 2023 & publicationPeriod: 2022-01-01 a 2023-12-31 \\ \hline
GT012 & T & cartas publicadas no primeiro trimestre deste ano & publicationPeriod: 1º trimestre do ano corrente \\ \hline
GT013 & T & cartas publicadas no segundo semestre do ano passado & publicationPeriod: 2º semestre do ano anterior \\ \hline
GT014 & T & cartas publicadas nesta semana & publicationPeriod: semana corrente \\ \hline
GT015 & T & cartas publicadas hoje & publicationPeriod: dia da execução \\ \hline
GT016 & C · T · A & cartas de pernambuco publicadas desde 2020 & state: "Pernambuco"; publicationPeriod: desde 2020 \\ \hline
GT017 & C · T · A & cartas da bahia publicadas nos últimos 6 meses & state: "Bahia"; publicationPeriod: últimos 6 meses \\ \hline
GT018 & C · T · A & cartas do AP produzidas no segundo semestre do ano passado & state: "Amapá"; creationPeriod: 2º semestre do ano anterior \\ \hline
GT019 & C · T · A & cartas do MA publicadas no ano passado & state: "Maranhão"; publicationPeriod: ano anterior, inteiro \\ \hline
GT020 & C · T & cartas da Paraíba publicadas esse ano & state: "Paraíba"; publicationPeriod: ano corrente \\ \hline
GT021 & C · T · A & cartas de roraima elaboradas hoje & state: "Roraima"; creationPeriod: dia da execução \\ \hline
GT022 & C · T · A & cartas de sergipe publicadas no primeiro trimestre deste ano & state: "Sergipe"; publicationPeriod: 1º trimestre do ano corrente \\ \hline
GT023 & C · T & cartas de Rondônia publicadas depois de 2022 & state: "Rondônia"; publicationPeriod: depois de 2022 \\ \hline
GT024 & C · T & cartas do Maranhão criadas nos últimos 3 meses & state: "Maranhão"; creationPeriod: últimos 3 meses \\ \hline
GT025 & C · T · A & cartas da bahia publicadas antes de 2020 & state: "Bahia"; publicationPeriod: antes de 2020 \\ \hline
GT026 & C · T · A & produtos do 3o cgeo publicados no ano passado & supplyArea: "3$^\circ$ Centro de Geoinformação"; publicationPeriod: ano anterior, inteiro \\ \hline
GT027 & C · T · A & produtos do 3 CGEO publicados hoje & supplyArea: "3$^\circ$ Centro de Geoinformação"; publicationPeriod: dia da execução \\ \hline
GT028 & C · T · A & produtos do 3º CGEO criados no primeiro trimestre deste ano & supplyArea: "3$^\circ$ Centro de Geoinformação"; creationPeriod: 1º trimestre do ano corrente \\ \hline
GT029 & C · T · A & produtos do quarto cgeo publicados depois de 2022 & supplyArea: "4$^\circ$ Centro de Geoinformação"; publicationPeriod: depois de 2022 \\ \hline
GT030 & C · T · A & produtos do 1º CGEO publicados no mês passado & supplyArea: "1$^\circ$ Centro de Geoinformação"; publicationPeriod: mês anterior, inteiro \\ \hline
GT031 & C · T · A & produtos do terceiro cgeo criados nos últimos 3 meses & supplyArea: "3$^\circ$ Centro de Geoinformação"; creationPeriod: últimos 3 meses \\ \hline
GT032 & C · T · A & produtos do 2 CGEO publicados antes de 2020 & supplyArea: "2$^\circ$ Centro de Geoinformação"; publicationPeriod: antes de 2020 \\ \hline
GT033 & C · T · A & produtos do 5 CGEO publicados na semana passada & supplyArea: "5$^\circ$ Centro de Geoinformação"; publicationPeriod: semana passada \\ \hline
GT034 & C · T & cartas na escala 1:5.000 publicadas depois de 2022 & scale: "1:5.000"; publicationPeriod: depois de 2022 \\ \hline
GT035 & C · T · A & cartas na escala 1:5000 publicadas esse ano & scale: "1:5.000"; publicationPeriod: ano corrente \\ \hline
GT036 & C · T · A & cartas na escala de 25 mil criadas no ano passado & scale: "1:25.000"; creationPeriod: ano anterior, inteiro \\ \hline
GT037 & C · T · A & cartas em 50k publicadas no segundo semestre do ano passado & scale: "1:50.000"; publicationPeriod: 2º semestre do ano anterior \\ \hline
GT038 & C · T · A & cartas na escala de 10 mil publicadas no mês passado & scale: "1:10.000"; publicationPeriod: mês anterior, inteiro \\ \hline
GT039 & C · T & cartas na escala 1:2.000 elaboradas hoje & scale: "1:2.000"; creationPeriod: dia da execução \\ \hline
GT040 & C · T · A & cartas na escala de 100 mil publicadas nesta semana & scale: "1:100.000"; publicationPeriod: semana corrente \\ \hline
GT041 & C · T · A & cartas na escala 1:5000 publicadas nos últimos 6 meses & scale: "1:5.000"; publicationPeriod: últimos 6 meses \\ \hline
\end{longtable}
\addtocounter{table}{-1}
}

\section{Consultas com ordena\c{c}\~ao geradas por \textit{templates} (G-O)}

Total desta fam\'ilia: \textbf{19 consultas}, listadas integralmente abaixo (a fun\c{c}\~ao geradora e o modelo de frase de cada consulta constam de \texttt{dataset.json}).

{\footnotesize
\begin{longtable}{|p{1.1cm}|p{1.3cm}|p{5.4cm}|p{6.3cm}|}
\hline
\textbf{ID} & \textbf{Cat.} & \textbf{Consulta} & \textbf{Leituras aceitas} \\ \hline
\endhead
GO001 & C · O · A & a carta mais recente do CE & state: "Ceará"; sortField: "publicationDate" ou "creationDate"; sortDirection: "DESC"; limit: 1 (opcional) \\ \hline
GO002 & C · O · A & a carta mais antiga do PA & state: "Pará"; sortField: "publicationDate" ou "creationDate"; sortDirection: "ASC"; limit: 1 (opcional) \\ \hline
GO003 & C · O & as cartas mais recentes do Ceará & state: "Ceará"; sortField: "publicationDate" ou "creationDate"; sortDirection: "DESC" \\ \hline
GO004 & C · O · A & as cartas mais antigas de mato grosso & state: "Mato Grosso"; sortField: "publicationDate" ou "creationDate"; sortDirection: "ASC" \\ \hline
GO005 & C · O · A & a primeira carta produzida de sergipe & state: "Sergipe"; sortField: "creationDate"; sortDirection: "ASC"; limit: 1 (opcional) \\ \hline
GO006 & C · O · A & a última carta publicada de GO & state: "Goiás"; sortField: "publicationDate"; sortDirection: "DESC"; limit: 1 (opcional) \\ \hline
GO007 & C · O · A & as 3 cartas mais antigas de sergipe & state: "Sergipe"; sortField: "publicationDate" ou "creationDate"; sortDirection: "ASC"; limit: 3 \\ \hline
GO008 & C · O · A & as 5 publicações mais recentes de MG & state: "Minas Gerais"; sortField: "publicationDate"; sortDirection: "DESC"; limit: 5 \\ \hline
GO009 & C · O · A & as 10 cartas criadas mais recentemente do PR & state: "Paraná"; sortField: "creationDate"; sortDirection: "DESC"; limit: 10 \\ \hline
GO010 & C · O · A & cartas do distrito federal em ordem cronológica de publicação & state: "Distrito Federal"; sortField: "publicationDate"; sortDirection: "ASC" \\ \hline
GO011 & C · O & as cartas temáticas mais recentes & productType: "Cartas Temáticas Não SCN"; sortField: "publicationDate" ou "creationDate"; sortDirection: "DESC" \\ \hline
GO012 & C · O & as 3 cartas CIRC mais antigas & productType: "CIRC"; sortField: "publicationDate" ou "creationDate"; sortDirection: "ASC"; limit: 3 \\ \hline
GO013 & C · O & a última ortoimagem publicada & productType: "SCN Carta Ortoimagem"; sortField: "publicationDate"; sortDirection: "DESC"; limit: 1 (opcional) \\ \hline
GO014 & C · O & a primeira carta topográfica vetorial elaborada & productType: "SCN Carta Topográfica Vetorial"; sortField: "creationDate"; sortDirection: "ASC"; limit: 1 (opcional) \\ \hline
GO015 & C · O & cartas topográficas da mais nova para a mais antiga & productType: "SCN Carta Topográfica Matricial"; sortField: "publicationDate" ou "creationDate"; sortDirection: "DESC" \\ \hline
GO016 & C · T · O & as cartas mais recentes publicadas entre 2022 e 2023 & publicationPeriod: 2022-01-01 a 2023-12-31; sortField: "publicationDate" ou "creationDate"; sortDirection: "DESC" \\ \hline
GO017 & C · T · O & a primeira carta criada na semana passada & creationPeriod: semana passada; sortField: "creationDate"; sortDirection: "ASC"; limit: 1 (opcional) \\ \hline
GO018 & C · T · O & as 5 cartas mais antigas publicadas nos últimos 5 anos & publicationPeriod: últimos 5 anos; sortField: "publicationDate" ou "creationDate"; sortDirection: "ASC"; limit: 5 \\ \hline
GO019 & C · T · O & a última carta publicada nos últimos 6 meses & publicationPeriod: últimos 6 meses; sortField: "publicationDate"; sortDirection: "DESC"; limit: 1 (opcional) \\ \hline
\end{longtable}
\addtocounter{table}{-1}
}

\section{Consultas amb\'iguas/informais adicionais (G-A)}

Total desta fam\'ilia: \textbf{24 consultas}, listadas integralmente abaixo (a fun\c{c}\~ao geradora e o modelo de frase de cada consulta constam de \texttt{dataset.json}).

{\footnotesize
\begin{longtable}{|p{1.1cm}|p{1.3cm}|p{5.4cm}|p{6.3cm}|}
\hline
\textbf{ID} & \textbf{Cat.} & \textbf{Consulta} & \textbf{Leituras aceitas} \\ \hline
\endhead
GA001 & S · A & mapas de pe & state: "Pernambuco" \\ \hline
GA002 & S · A & mapas do pr & state: "Paraná" \\ \hline
GA003 & S · A & mapas de sc & state: "Santa Catarina" \\ \hline
GA004 & S · A & mapas de go & state: "Goiás" \\ \hline
GA005 & S · A & mapas do ce & state: "Ceará" \\ \hline
GA006 & S · A & mapas do pa & state: "Pará" \\ \hline
GA007 & S · A & mapas de goias & state: "Goiás" \\ \hline
GA008 & S · A & mapas de rondonia & state: "Rondônia" \\ \hline
GA009 & S · A & mapas do piaui & state: "Piauí" \\ \hline
GA010 & S · A & mapas do amapa & state: "Amapá" \\ \hline
GA011 & S · A & mapas do maranhao & state: "Maranhão" \\ \hline
GA012 & S · A & mapas da paraiba & state: "Paraíba" \\ \hline
GA013 & S · A & mapas de cuiaba & city: "Cuiabá" \\ \hline
GA014 & S · A & mapas de goiania & city: "Goiânia" \\ \hline
GA015 & S · A & mapas de belem & city: "Belém" \\ \hline
GA016 & S · A & mapas de brasilia & city: "Brasília" \\ \hline
GA017 & C · A & cartas do rio em 50k & state: "Rio de Janeiro"; scale: "1:50.000"; aceita-se também city no lugar de state \\ \hline
GA018 & C · A & mapas detalhados do Amazonas & scale: "1:25.000"; state: "Amazonas" \\ \hline
GA019 & S · A & cartas em pequena escala do sul & scale: "1:250.000" \\ \hline
GA020 & S · A & produtos em média escala & scale: "1:50.000" ou "1:100.000" \\ \hline
GA021 & C · A & cartas de Goiás em 25k & state: "Goiás"; scale: "1:25.000" \\ \hline
GA022 & C · A & cartas do Tocantins em 50k & state: "Tocantins"; scale: "1:50.000" \\ \hline
GA023 & C · A & cartas de Mato Grosso do Sul em 100k & state: "Mato Grosso do Sul"; scale: "1:100.000" \\ \hline
GA024 & C · A & cartas de Pernambuco em 250k & state: "Pernambuco"; scale: "1:250.000" \\ \hline
\end{longtable}
\addtocounter{table}{-1}
}

\section{Consultas amplas, com muitos produtos (G-P)}

Total desta fam\'ilia: \textbf{13 consultas}, listadas integralmente abaixo (a fun\c{c}\~ao geradora e o modelo de frase de cada consulta constam de \texttt{dataset.json}).

{\footnotesize
\begin{longtable}{|p{1.1cm}|p{1.3cm}|p{5.4cm}|p{6.3cm}|}
\hline
\textbf{ID} & \textbf{Cat.} & \textbf{Consulta} & \textbf{Leituras aceitas} \\ \hline
\endhead
GP001 & S & todas as cartas de Rondônia & state: "Rondônia" \\ \hline
GP002 & S & todas as cartas do Tocantins & state: "Tocantins" \\ \hline
GP003 & S & todas as cartas do Piauí & state: "Piauí" \\ \hline
GP004 & S & todas as cartas de Mato Grosso & state: "Mato Grosso" \\ \hline
GP005 & S & todas as cartas do Acre & state: "Acre" \\ \hline
GP006 & S & todos os produtos do 1$^\circ$ Centro de Geoinformação & supplyArea: "1$^\circ$ Centro de Geoinformação" \\ \hline
GP007 & S & todos os produtos do 2$^\circ$ Centro de Geoinformação & supplyArea: "2$^\circ$ Centro de Geoinformação" \\ \hline
GP008 & S & todos os produtos do 3$^\circ$ Centro de Geoinformação & supplyArea: "3$^\circ$ Centro de Geoinformação" \\ \hline
GP009 & S & todos os produtos do 4$^\circ$ Centro de Geoinformação & supplyArea: "4$^\circ$ Centro de Geoinformação" \\ \hline
GP010 & S & todos os produtos do 5$^\circ$ Centro de Geoinformação & supplyArea: "5$^\circ$ Centro de Geoinformação" \\ \hline
GP011 & C & cartas na escala 1:2.000 do Rio Grande do Norte & scale: "1:2.000"; state: "Rio Grande do Norte" \\ \hline
GP012 & C & cartas na escala 1:5.000 de Alagoas & scale: "1:5.000"; state: "Alagoas" \\ \hline
GP013 & C & cartas na escala 1:250.000 do Rio Grande do Sul & scale: "1:250.000"; state: "Rio Grande do Sul" \\ \hline
\end{longtable}
\addtocounter{table}{-1}
}

\section{Consultas de fronteira do dom\'inio (G-F)}

Total desta fam\'ilia: \textbf{7 consultas}, listadas integralmente abaixo (a fun\c{c}\~ao geradora e o modelo de frase de cada consulta constam de \texttt{dataset.json}).

{\footnotesize
\begin{longtable}{|p{1.1cm}|p{1.3cm}|p{5.4cm}|p{6.3cm}|}
\hline
\textbf{ID} & \textbf{Cat.} & \textbf{Consulta} & \textbf{Leituras aceitas} \\ \hline
\endhead
GF001 & F & me ajuda & não chamar a ferramenta (sem intenção de busca) \\ \hline
GF002 & F & cartas de Marte & não chamar a ferramenta (fora da Terra) \\ \hline
GF003 & F & cartas da Argentina & não chamar a ferramenta (fora do território brasileiro) \\ \hline
GF004 & F & qual a previsão do tempo para amanhã em Manaus? & não chamar a ferramenta (outro assunto) \\ \hline
GF005 & F & quanto custa uma carta topográfica? & não chamar a ferramenta (pergunta de preço, não de busca) \\ \hline
GF006 & F & mapa do tesouro pirata & não chamar a ferramenta (ficção) \\ \hline
GF007 & S & cartas do estado 42 & não chamar a ferramenta (observacional, P6: valor impossível de representar) \\ \hline
\end{longtable}
\addtocounter{table}{-1}
}


\shorthandon{"}

\section{Requisitos exercitados pelas consultas}
\label{apx:dataset-requisitos}

O Quadro~\ref{quadro:requisitos-por-consulta} relaciona os grupos de consultas aos requisitos da Seção~\ref{sec:requisitos}. Todas as consultas exercitam a tradução (RF2) e o registro dos traços de execução (RF5), pois cada rodada do \textit{pipeline} de avaliação grava um registro por consulta. A troca do modelo entre rodadas, por parâmetro do \textit{pipeline}, exercita a substituição do modelo por configuração (RNF6). A rota REST (RF1), a busca no banco (RF3) e a seleção do modelo por requisição (RF4) ficam fora das métricas, que avaliam a tradução (Quadro~\ref{quadro:requisitos-resumo}).

\begin{quadro}[htbp!]
\centering
\caption{Grupos de consultas do \textit{dataset} e aspectos da tradução que exercitam}
\label{quadro:requisitos-por-consulta}
\small
\begin{tabular}{|p{3.9cm}|p{7.9cm}|c|}
\hline
\textbf{Grupo de consultas} & \textbf{O que exige do modelo} & \textbf{Requisitos} \\
\hline
Todas (310) & Preencher os campos da ferramenta sustentados pela consulta, na forma canônica, sem acrescentar campos & RF2, RF5, RNF6 \\
\hline
Com código MI/INOM (M) & Preservar o código em \texttt{keyword}, normalizando grafias compactadas do INOM & RF2 \\
\hline
Com tempo relativo (T) & Resolver a expressão em datas ISO a partir da data de referência informada e escolher entre período de publicação e de criação & RF2 \\
\hline
Com ordenação (O) & Preencher \texttt{sortField}, \texttt{sortDirection} e, quando pedido, \texttt{limit}, sem ordenar quando a consulta não pede & RF2 \\
\hline
Ambíguas/informais (A) & Normalizar siglas, grafias sem acento, apelidos e expressões qualitativas sem dicionário de normalização & RF2 \\
\hline
Fora do domínio (F) & Reconhecer que o catálogo não atende ao pedido e não chamar a ferramenta & RF2 \\
\hline
Observacionais (4) & Comportamento diante de valores impossíveis ou de pedidos de resposta discutível (reportado, não pontuado) & RF2, RF5 \\
\hline
\end{tabular}
\fonte{Elaborado pelos autores. As contagens por categoria estão na Tabela~\ref{tab:dataset-categorias}.}
\end{quadro}

\section{Auditoria do \textit{dataset}}
\label{apx:auditoria}

A auditoria segue o procedimento do manual de anotação (Seção~\ref{apx:manual}, subseção ``Procedimento de auditoria'') e é reproduzível a partir do repositório. Os artefatos ficam em \texttt{data/auditoria/}. Eles guardam o resultado das checagens automáticas e das medidas de concordância, uma linha por consulta com as evidências e as divergências, a anotação às cegas com a justificativa de cada consulta, as decisões de adjudicação com o gabarito anterior a cada uma e a amostra estratificada de 40 consultas (semente 42) conferida por um dos autores. O comando \texttt{pfc-auditar} (Seção~\ref{apx:comandos}) refaz a auditoria e regenera as Tabelas~\ref{tab:auditoria} e~\ref{tab:auditoria-campos} e o Quadro~\ref{quadro:adjudicacao}.

A Tabela~\ref{tab:auditoria-campos} detalha, por campo, a concordância entre a leitura preferencial do gabarito e a da anotação independente. O Quadro~\ref{quadro:adjudicacao} lista todas as decisões de adjudicação. Quatro consultas tiveram leitura preferencial divergente (Tabela~\ref{tab:auditoria}). Em P13, P21 e P22, a divergência decorre da regra que mantém, na camada P, a anotação original da equipe do 1º CGEO como preferencial. Em N33, o gabarito inclui na leitura preferencial o campo opcional \texttt{project}, que o anotador deixou de fora. Nas quatro, a leitura do anotador está entre as leituras aceitas.

\begin{table}[htbp!]
\centering
\caption{Concordância por campo entre o gabarito e a anotação independente (leitura preferencial, após a adjudicação)}
\label{tab:auditoria-campos}
\small
\begin{tabular}{|l|c|c|c|c|c|}
\hline
\textbf{Campo} & \textbf{Gabarito} & \textbf{Anotador} & \textbf{Ambos} & \textbf{Kappa} & \textbf{Valor igual} \\
\hline
\texttt{keyword} & 62 & 62 & 62 & 1,000 & 62/62 \\
\texttt{scale} & 84 & 84 & 84 & 1,000 & 83/84 \\
\texttt{productType} & 44 & 44 & 44 & 1,000 & 44/44 \\
\texttt{state} & 130 & 130 & 130 & 1,000 & 130/130 \\
\texttt{city} & 27 & 27 & 27 & 1,000 & 27/27 \\
\texttt{supplyArea} & 40 & 40 & 40 & 1,000 & 40/40 \\
\texttt{project} & 11 & 10 & 10 & 0,951 & 10/10 \\
\texttt{publicationPeriod} & 51 & 51 & 51 & 1,000 & 51/51 \\
\texttt{creationPeriod} & 10 & 10 & 10 & 1,000 & 10/10 \\
\texttt{sortField} & 29 & 29 & 29 & 1,000 & 27/29 \\
\texttt{sortDirection} & 29 & 29 & 29 & 1,000 & 29/29 \\
\texttt{limit} & 21 & 21 & 21 & 1,000 & 21/21 \\
\hline
\end{tabular}
\fonte{Elaborada pelos autores. Gabarito, Anotador e Ambos: consultas em que o campo está na leitura preferencial de cada um. Kappa de Cohen sobre a presença do campo. Valor igual: entre as consultas em que ambos o anotam; períodos comparados pelos limites, textos livres após normalização.}
\end{table}


\begin{quadro}[htbp!]
\centering
\caption{Decisões de adjudicação da auditoria do \textit{dataset}}
\label{quadro:adjudicacao}
\footnotesize
\begin{tabular}{|p{2.3cm}|p{4.9cm}|p{1.5cm}|p{5.3cm}|}
\hline
\textbf{Consultas} & \textbf{Divergência} & \textbf{Decisão} & \textbf{Justificativa} \\
\hline
N36, N37, GF001, GF002, GF003, GF004, GF005, GF006\newline{\scriptsize\itshape anotação independente} & o gabarito classificava a consulta como observacional; o anotador a tratou como fora do domínio, com gabarito determinado 'não chamar' & corrigido & P4 e §3 do manual: consulta sem intenção de busca no acervo brasileiro tem comportamento correto determinado (não chamar); P6 reserva 'observacional' às consultas cujo comportamento correto não é determinável. A consulta passa às métricas principais, na categoria F (fora do domínio). \\ \hline
N33\newline{\scriptsize\itshape anotação independente} & o gabarito exigia project = Mapeamento Sistemático; o anotador anotou só state = São Paulo e listou project como leitura alternativa, pois 'cartografia sistemática' não é nome nem apelido informado na ferramenta & ampliado & As duas leituras são razoáveis (P3): a expressão designa o programa pelo termo distintivo do nome, mas não consta da descrição da ferramenta. project passa a campo opcional; a regra de project do manual passa a registrar o caso. \\ \hline
GF007\newline{\scriptsize\itshape anotação independente} & ambos classificam a consulta como observacional; o gabarito indica 'não chamar' e o anotador, 'chamar sem parâmetros' & mantido & Consulta observacional (P6, valor impossível de representar): o comportamento é reportado, mas não pontuado, de modo que a divergência não afeta nenhuma métrica. Mantém-se a indicação 'não chamar', coerente com N38 (escala 1:10.000.000). \\ \hline
P14\newline{\scriptsize\itshape consistência} & 'detalhado' é anotado como 1:25.000 em todo o dataset, mas P14 aceita 1:25.000 ou 1:50.000 & mantido & As escalas explícitas '(25k ou 50k)' prevalecem sobre o qualificativo 'detalhado' (regra de scale: duas escalas explícitas, aceita-se qualquer uma). \\ \hline
P13\newline{\scriptsize\itshape anotação independente (medida estrita)} & leitura preferencial diferente dentro do conjunto aceito: scale 1:25.000 no gabarito, 1:50.000 no anotador (ambas aceitas por 'maior que 100k') & mantido & Na camada P, a anotação original da equipe do 1º CGEO é mantida como leitura preferencial (manual, §4); a escolha da preferencial não afeta a pontuação. \\ \hline
P21, P22\newline{\scriptsize\itshape anotação independente (medida estrita)} & leitura preferencial diferente dentro do conjunto aceito: sortField creationDate no gabarito, publicationDate no anotador (ambos aceitos, sem pista ligada à ordenação) & mantido & Na camada P, a anotação original da equipe do 1º CGEO é mantida como leitura preferencial (manual, §4); a escolha da preferencial não afeta a pontuação. \\ \hline
GO016\newline{\scriptsize\itshape anotação independente (medida estrita)} & o gabarito aceita sortField publicationDate ou creationDate; o anotador aceitou só publicationDate & mantido & Em 'as cartas mais recentes publicadas entre 2022 e 2023', o verbo qualifica o filtro de período, não a ordenação; sem pista ligada à ordenação, aceitam-se os dois campos (regra de ordenação). É o mesmo raciocínio do anotador em P22. \\ \hline
\end{tabular}
\fonte{Elaborado pelo autor a partir de \texttt{data/auditoria/adjudicacao.json}.}
\end{quadro}


\clearpage
\section{Estrutura, composição e concordância dos conjuntos}
\label{apx:dataset-tabelas}

Esta seção reúne as tabelas que descrevem as 310 consultas e os dois lotes (Seção~\ref{sec:dataset}). Além dos campos da Tabela~\ref{tab:estrutura-dataset}, cada entrada traz os campos \texttt{origem} (P, N ou G), \texttt{familia} (P, N, GS \ldots GF) e \texttt{trocas}, com as leituras estruturais alternativas que o campo \texttt{alternativas} expande. Os valores alternativos e os campos opcionais ficam marcados no próprio \texttt{esperado}.

\begin{table}[htbp!]
\centering
\caption{Estrutura de cada entrada do \textit{dataset} de avaliação}
\label{tab:estrutura-dataset}
\small
\begin{tabular}{|l|p{10.5cm}|}
\hline
\textbf{Campo} & \textbf{Descrição} \\
\hline
\texttt{id} & Identificador (P01--P22, N01--N40, GS001 \ldots GF007) \\ \hline
\texttt{fonte} & Rastreabilidade: arquivo e linha de origem (camada P), autores (camada N) ou função geradora, modelo de frase e semente (camada G) \\ \hline
\texttt{categorias} & Uma ou mais das categorias S, C, M, T, O, A; ou F (fora do domínio), que não se combina com as demais \\ \hline
\texttt{consulta} & Texto em linguagem natural \\ \hline
\texttt{esperado} & Leitura preferencial do gabarito: parâmetros de busca na forma canônica, com períodos relativos expressos por regras nomeadas \\ \hline
\texttt{alternativas} & Demais leituras aceitas, quando a consulta admite mais de uma (Seção~\ref{sec:manual-anotacao}) \\ \hline
\texttt{espera\_tool\_call} & Falso quando o comportamento correto é não chamar a ferramenta \\ \hline
\texttt{observacional} & Verdadeiro quando o caso fica fora das métricas principais \\ \hline
\texttt{notas} & Justificativa das leituras alternativas e observações do anotador \\ \hline
\end{tabular}%
\fonte{Elaborada pelos autores.}
\end{table}

\begin{table}[htbp!]
\centering
\caption{Consultas por categoria e por camada (métricas principais; uma consulta pode ter mais de uma categoria)}
\label{tab:dataset-categorias}
\small
\begin{tabular}{|l|c|c|c|c|}
\hline
\textbf{Categoria} & \textbf{P} & \textbf{N} & \textbf{G} & \textbf{Total} \\
\hline
Simples (S) & 0 & 11 & 89 & 100 \\
Compostas (C) & 22 & 21 & 137 & 180 \\
Código MI/INOM (M) & 10 & 5 & 42 & 57 \\
Tempo relativo (T) & 6 & 9 & 45 & 60 \\
Ordenação (O) & 4 & 6 & 19 & 29 \\
Ambíguas/informais (A) & 16 & 19 & 162 & 197 \\
Fora do domínio (F) & 0 & 2 & 6 & 8 \\
\hline
\end{tabular}
\fonte{Elaborada pelo autor.}
\end{table}


\begin{table}[htbp!]
\centering
\caption{Frequência de cada campo na leitura preferencial do gabarito (métricas principais)}
\label{tab:dataset-campos}
\small
\begin{tabular}{|l|c||l|c|}
\hline
\textbf{Campo} & \textbf{Ocorr.} & \textbf{Campo} & \textbf{Ocorr.} \\
\hline
\texttt{keyword} & 62 & \texttt{project} & 11 \\
\texttt{scale} & 84 & \texttt{publicationPeriod} & 51 \\
\texttt{productType} & 44 & \texttt{creationPeriod} & 10 \\
\texttt{state} & 130 & \texttt{sortField} & 29 \\
\texttt{city} & 27 & \texttt{sortDirection} & 29 \\
\texttt{supplyArea} & 40 & \texttt{limit} & 21 \\
\hline
\end{tabular}
\fonte{Elaborada pelos autores.}
\end{table}


O lote 1 (validação) e o lote 2 (teste) foram construídos com a mesma receita (Seção~\ref{sec:lote-validacao}). As Tabelas~\ref{tab:lote_composicao} e~\ref{tab:lote2_composicao} dão a composição de cada lote, e as Tabelas~\ref{tab:lote_concordancia} e~\ref{tab:lote2_concordancia}, a concordância da anotação às cegas.

\begin{table}[htbp!]
\centering
\caption{Lote de validação: do alvo ao lote final, por família}
\label{tab:lote_composicao}
\footnotesize
\begin{tabular}{|l|l|c|c|c|c|}
\hline
\textbf{Família} & \textbf{Testa} & \textbf{Alvos} & \textbf{Checagens} & \textbf{Anotação às cegas} & \textbf{Lote final} \\
\hline
VS & Simples (um critério) & 280 & $-$3 & 0 & 277 \\
VC & Compostas & 360 & 0 & $-$1 & 359 \\
VM & Códigos MI/INOM & 160 & 0 & 0 & 160 \\
VT & Tempo & 260 & $-$1 & 0 & 259 \\
VO & Ordenação & 200 & $-$1 & 0 & 199 \\
VA & Leituras múltiplas & 200 & $-$7 & 0 & 193 \\
VE & Subespecificadas & 150 & $-$4 & 0 & 146 \\
VF & Fora do domínio & 340 & $-$4 & 0 & 336 \\
\hline
\textbf{Total} & & 1.950 & $-$20 & $-$1 & \textbf{1.929} \\
\hline
\end{tabular}
\fonte{Elaborada pelos autores. Checagens: descartadas nas checagens automáticas e na deduplicação. Anotação às cegas: descartadas por incompatibilidade entre o gabarito por construção e a anotação A.}
\end{table}


\begin{table}[htbp!]
\centering
\caption{Concordância na anotação às cegas do lote de validação}
\label{tab:lote_concordancia}
\footnotesize
\begin{tabular}{|l|c|c|}
\hline
\textbf{Medida} & \textbf{Construção $\times$ A} (n=1.950) & \textbf{A $\times$ B} (n=390) \\
\hline
Decisão (buscar / buscar ou esclarecer / não buscar), kappa & 1,000 & 1,000 \\
Leituras compatíveis nos dois sentidos & 99,8\% & 100,0\% \\
Leitura preferencial idêntica & 99,8\% & 100,0\% \\
Presença dos campos, kappa & 0,999 & 1,000 \\
Valor igual quando ambos preenchem & 100,0\% & 100,0\% \\
Detecção de leituras múltiplas, kappa & 0,967 & 1,000 \\
\hline
\end{tabular}
\fonte{Elaborada pelos autores. Medidas de \texttt{audit.concordancia}, as mesmas da auditoria das 310 consultas, sobre todas as consultas redigidas (antes do filtro). A: anotação de todas as consultas; B: segunda anotação independente de uma amostra aleatória de 20\%.}
\end{table}


\begin{table}[htbp!]
\centering
\caption{Do alvo ao lote final, por família, no lote 2 (teste)}
\label{tab:lote2_composicao}
\footnotesize
\begin{tabular}{|l|l|c|c|c|c|}
\hline
\textbf{Família} & \textbf{Testa} & \textbf{Alvos} & \textbf{Checagens} & \textbf{Anotação às cegas} & \textbf{Lote final} \\
\hline
VS & Simples (um critério) & 140 & $-$5 & $-$1 & 134 \\
VC & Compostas & 180 & 0 & $-$3 & 177 \\
VM & Códigos MI/INOM & 80 & $-$1 & 0 & 79 \\
VT & Tempo & 130 & 0 & 0 & 130 \\
VO & Ordenação & 100 & $-$1 & 0 & 99 \\
VA & Leituras múltiplas & 100 & $-$2 & 0 & 98 \\
VE & Subespecificadas & 75 & $-$10 & 0 & 65 \\
VF & Fora do domínio & 170 & $-$7 & 0 & 163 \\
\hline
\textbf{Total} & & 975 & $-$26 & $-$4 & \textbf{945} \\
\hline
\end{tabular}
\fonte{Elaborada pelos autores. Checagens: descartadas nas checagens automáticas e na deduplicação. Anotação às cegas: descartadas por incompatibilidade entre o gabarito por construção e a anotação A.}
\end{table}


\begin{table}[htbp!]
\centering
\caption{Concordância na anotação às cegas do lote 2 (teste)}
\label{tab:lote2_concordancia}
\footnotesize
\ajustartabela{%
\begin{tabular}{|l|c|c|}
\hline
\textbf{Medida} & \textbf{Construção $\times$ A} (n=975) & \textbf{A $\times$ B} (n=195) \\
\hline
Decisão (buscar / buscar ou esclarecer / não buscar), kappa & 1,000 & 1,000 \\
Leituras compatíveis nos dois sentidos & 99,6\% & 100,0\% \\
Leitura preferencial idêntica & 99,5\% & 100,0\% \\
Presença dos campos, kappa & 0,998 & 1,000 \\
Valor igual quando ambos preenchem & 100,0\% & 100,0\% \\
Detecção de leituras múltiplas, kappa & 0,961 & 1,000 \\
\hline
\end{tabular}}
\fonte{Elaborada pelos autores. Medidas de \texttt{audit.concordancia}, as mesmas da auditoria das 310 consultas, sobre todas as consultas redigidas (antes do filtro). A: anotação de todas as consultas; B: segunda anotação independente de uma amostra aleatória de 20\%.}
\end{table}


% ----------------------------------------------------------
% Apêndice B — Especificação dos Ambientes de Execução
% ----------------------------------------------------------
\chapter{Especificação dos Ambientes de Execução}
\label{apx:specs}

\begin{sloppypar}
Este apêndice documenta o \textit{hardware} e o \textit{software} dos ambientes de execução da Seção~\ref{sec:ambiente}. Em todos eles, cada rodada grava um manifesto (\texttt{results/<modelo>/manifesto.json}; Seção~\ref{sec:reprodutibilidade}).
\end{sloppypar}

\section{Estação de referência}
\label{apx:specs-estacao}

A estação de referência tem a seguinte configuração:
\begin{itemize}
  \item \textbf{CPU}. Intel Core i7-12700H (12.ª geração, 14 núcleos / 20 \textit{threads}, base 2,3~GHz).
  \item \textbf{RAM}. 32~GB DDR4 a 3200~MHz.
  \item \textbf{GPU dedicada}. NVIDIA GeForce RTX 3050 Laptop GPU com 4~GB de VRAM.
  \item \textbf{Disco}. SSD NVMe de 512~GB.
  \item \textbf{Sistema operacional}. Microsoft Windows 11 Pro (versão 10.0.26200, 64 bits).
  \item \textbf{Software}. Python 3.14.4; Ollama 0.34.0; LangChain 1.x (\texttt{langchain-core} 1.6.4, \texttt{langchain-ollama} 1.1.0, cliente \texttt{ollama} 0.6.2); PostgreSQL 16.4 com PostGIS 3.4 em contêiner Docker, usado apenas na demonstração ponta a ponta.
\end{itemize}

Os dados foram coletados no próprio equipamento, em setembro de 2026, pelos comandos PowerShell da Seção~\ref{apx:specs-comandos}, com o mesmo \textit{driver} NVIDIA das rodadas da estação. A Seção~\ref{apx:specs-resultado} transcreve a saída de forma condensada. A imagem Docker do banco da demonstração, \texttt{postgis/postgis:16-3.4}, está definida no \texttt{docker-compose.yml} do repositório. O \texttt{psql} listado é o cliente instalado na estação.

\subsection{Comandos PowerShell de coleta}
\label{apx:specs-comandos}

{\footnotesize
\begin{verbatim}
# CPU
Get-CimInstance Win32_Processor | Select-Object Name, NumberOfCores,
  NumberOfLogicalProcessors, MaxClockSpeed

# Memoria RAM (total em GB e detalhes por pente)
(Get-CimInstance Win32_PhysicalMemory |
  Measure-Object -Property Capacity -Sum).Sum / 1GB
Get-CimInstance Win32_PhysicalMemory | Select-Object Manufacturer,
  Capacity, Speed, ConfiguredClockSpeed, MemoryType, DeviceLocator

# GPU (geral)
Get-CimInstance Win32_VideoController | Select-Object Name,
  AdapterRAM, DriverVersion, VideoProcessor

# VRAM real (se GPU NVIDIA)
nvidia-smi --query-gpu=name,memory.total,driver_version --format=csv

# Sistema operacional
Get-CimInstance Win32_OperatingSystem | Select-Object Caption,
  Version, BuildNumber, OSArchitecture

# Disco principal
Get-PhysicalDisk | Select-Object FriendlyName, MediaType, Size, BusType

# Versoes de software
python --version
ollama --version
psql --version
docker --version
\end{verbatim}
}

\subsection{Saída coletada}
\label{apx:specs-resultado}

{\scriptsize
\begin{verbatim}
# CPU
Name                                 NumberOfCores  NumberOfLogicalProcessors  MaxClockSpeed
12th Gen Intel(R) Core(TM) i7-12700H            14                         20           2300

# RAM
RAM total (GB): 32
Pente 1: Kingston, 16 GB DDR4, 3200 MHz (Controller0-ChannelA-DIMM0)
Pente 2: Kingston, 16 GB DDR4, 3200 MHz (Controller1-ChannelA-DIMM0)

# GPU
Name: NVIDIA GeForce RTX 3050 Laptop GPU
  AdapterRAM 4 293 918 720; DriverVersion 32.0.16.1664
  VideoProcessor NVIDIA GeForce RTX 3050 Laptop GPU
Name: Intel(R) Iris(R) Xe Graphics
  AdapterRAM 2 147 479 552; DriverVersion 32.0.101.6790
  VideoProcessor Intel(R) Iris(R) Xe Graphics Family

# nvidia-smi
NVIDIA GeForce RTX 3050 Laptop GPU, 4096 MiB, 616.64

# Sistema operacional
Microsoft Windows 11 Pro, versao 10.0.26200, Build 26200, 64 bits

# Disco principal
INTEL SSDPEKNU512GZ, SSD NVMe, 512 GB

# Software
Python 3.14.4
ollama version is 0.34.0
psql (PostgreSQL) 18.3
Docker version 29.6.1, build 8900f1d
\end{verbatim}
}

A Tabela~\ref{tab:ambiente_estacao} transcreve dos manifestos o ambiente das duas rodadas de validação na estação, com o Qwen 3 4B e o Gemma 4 E2B (Seção~\ref{sec:validacao-local}). As rodadas feitas depois na estação, para o teste A/B, têm as condições de \textit{hardware} nos próprios manifestos, resumidas em \texttt{results/\allowbreak ab/\allowbreak ab.md}.

\tabelagerada{tab_ambiente_estacao}{ambiente da validação na estação}

\section{Rodadas completas em GPU de nuvem}
\label{apx:specs-colab}

As rodadas completas dos quatro modelos locais foram feitas numa máquina virtual do Google Colab com GPU NVIDIA T4 (16~GB nominais de VRAM, dos quais o \texttt{nvidia-smi} informa 15.360~MiB; Tabela~\ref{tab:ambiente}), pelo caderno \texttt{notebooks/avaliacao\_colab.ipynb} do repositório. O caderno instala o Ollama pelo instalador oficial e baixa os modelos pelas \textit{tags} \texttt{qwen3:4b-instruct-\allowbreak 2507-q4\_K\_M}, \texttt{gemma4:\allowbreak e4b-it-qat}, \texttt{gemma4:\allowbreak e2b-it-qat} e \texttt{mistral-nemo:12b}. As \textit{tags} do Qwen 3 4B, do Gemma 4 E4B e do Gemma 4 E2B são as mesmas usadas na estação de referência, com o mesmo \textit{digest} registrado nos manifestos. O caderno recebe o pacote do projeto (\texttt{pfc\_busca\_colab.zip}, gerado por \texttt{scripts/\allowbreak empacotar\_colab.py}), confere o \textit{hash} SHA-256 do conjunto de consultas gravado no próprio caderno e recusa-se a rodar com outro pacote. Depois, executa \texttt{pfc-avaliar} com três repetições por modelo, data de referência fixa (14 de setembro de 2026) e sem execução de SQL, pois a busca no banco não entra nas métricas.

No mesmo ambiente, o caderno \texttt{notebooks/linha\_de\_base\_colab.ipynb} executou, com uma repetição e a mesma data de referência, a Saída Estruturada v1 nos quatro modelos locais e o método do protótipo no \texttt{phi4:14b}. Os lotes 1 (validação) e 2 (teste) também rodaram nesse ambiente, com o Gemma 4 E4B e uma repetição. O caderno \texttt{notebooks/\allowbreak lote\_validacao\_colab.ipynb} executou as configurações v1 e v2 no lote 1 e as duas v2 nas 310. O \texttt{notebooks/\allowbreak lote2\_v3\_colab.ipynb} conferiu o \textit{hash} da v3 congelada antes de rodar e executou as configurações do plano de análise no lote 2 e a v3 e o controle nas 310. O \texttt{notebooks/\allowbreak lote2\_v3\_modelos\_colab.ipynb} executou a extensão ao Gemma 4 E2B e ao Qwen 3 4B no lote 2. Nenhum serviço de LLM em nuvem é usado nesse ambiente. O Ollama executa dentro da máquina virtual, com os pesos baixados nela, e só a GPU vem da nuvem.

\tabelagerada{tab_ambiente}{ambiente das rodadas completas}

\section{Referência em nuvem (Groq)}
\label{apx:specs-groq}

A referência em nuvem usa a API do Groq, compatível com a da OpenAI, pela classe \texttt{ChatOpenAI} do pacote \texttt{langchain-openai} (versão 1.6.6), com \texttt{langchain-core} 1.6.4, a partir da estação de referência. O \textit{prompt} de sistema, a definição da ferramenta, a temperatura (zero) e o limite de 1024 \textit{tokens} de saída são os mesmos das rodadas locais. O esforço de raciocínio foi fixado no menor valor aceito por cada modelo, \res{groq}{geral}{esforcos}. O valor \texttt{none} desliga o raciocínio. Cada modelo executou as \res{groq}{geral}{nexecutadas} consultas uma vez, com data de referência de 24 de setembro de 2026 e intervalo mínimo de 11~s entre chamadas.

Quando o provedor sinaliza o esgotamento do limite de taxa do plano gratuito, a execução aguarda o tempo indicado e retoma a mesma consulta, sem contá-la como erro. As falhas transitórias de conexão ou do servidor são repetidas até três vezes antes de contarem como erro de infraestrutura, e só a latência da tentativa respondida é registrada. O servidor valida os argumentos contra o \textit{schema}, e uma chamada de ferramenta rejeitada por argumentos inválidos conta como erro do modelo. Quando os argumentos podem ser recuperados da mensagem de rejeição, eles são avaliados campo a campo, como na execução local. A linha de base com Saída Estruturada usa o mesmo cliente, sem ferramenta, com o modo JSON do provedor (\texttt{response\_format} do tipo \texttt{json\_object}), a mesma data de referência e o mesmo esforço de raciocínio de cada modelo.

% ----------------------------------------------------------
% Apêndice C — Repositório de Código e Reprodução
% ----------------------------------------------------------
\chapter{Repositório de Código e Reprodução dos Experimentos}
\label{apendice:codigo}

O código-fonte, o \textit{dataset}, o manual de anotação, os registros de auditoria, os registros de execução de todas as rodadas e as fontes deste texto estão no repositório público do projeto, \url{https://github.com/vhenequim/bdgex-busca-tool-calling}. O código, o \textit{dataset} e os artefatos de avaliação estão sob a licença MIT. As partes derivadas do protótipo do 1º CGEO mantêm a licença MIT original, como registra o arquivo \texttt{THIRD\_PARTY\_NOTICES.md}. O texto deste trabalho segue os termos de uso do IME impressos no verso da folha de rosto. A estrutura principal do repositório aparece abaixo.

{\footnotesize
\begin{verbatim}
bdgex-busca-tool-calling/
  pyproject.toml         dependencias e comandos pfc e pfc-*
  src/pfc_busca/
    abordagens.py        registro unico das abordagens (familia, versao,
                         prefixo, tradutor e hashes do manifesto)
    cli.py               comando pfc (avaliar, relatorio, conferir, lote,
                         lote2, ab, abordagens, api, entre outros)
    schema.py            definicao da ferramenta buscar_catalogo (12 campos)
    prompts.py           prompt de sistema (com a data de referencia)
    agent.py             traducao via ChatOllama.bind_tools() (modelos locais)
    agent_groq.py        traducao via ChatOpenAI (referencia em nuvem)
    agent_estruturado.py linhas de base com Saida Estruturada e metodo
                         do prototipo
    prototipo_termos.py  dicionario do prototipo (so na linha de base)
    v2.py                especificacao v2 (descricoes e recusar_consulta)
    v3.py                v3 (laco de chamadas, retorno, recusa contestada
                         e controle com Saida Estruturada)
    ferramentas.py       ferramentas auxiliares e validador da v3
    dados/               municipios do IBGE e indice de folhas do BDGEx
                         (o indice nao e versionado)
    tools.py, db.py      busca no PostgreSQL/PostGIS (psycopg)
    pipeline.py          ciclo completo (traducao, busca e resposta)
    api.py               rotas POST /api/search e GET /api/health (FastAPI)
    evaluation/
      manual_cases.py    camadas P e N, com rastreabilidade e notas
      gabarito.py        leituras multiplas (valor alternativo, opcional, trocas)
      relative_time.py   regras de tempo relativo (leituras aceitas)
      dataset_builder.py pfc-dataset, constroi e valida data/dataset.json
      audit.py           pfc-auditar, auditoria do dataset
      run_evaluation.py  pfc-avaliar, execucao, manifesto, repontuacao
      metrics.py         acuracia, P/R/F1 por campo, latencia
      report.py          pfc-relatorio, tabelas, figuras, Wilson, McNemar
      conferencia.py     pfc-conferir, conferencia das rodadas
      comparacao.py      Tool Calling x Saida Estruturada (tabela e macros)
      lote.py            analise do lote 1 (tabelas e macros)
      lote2.py           analise do lote 2 (plano de analise da v3)
      ab.py              pfc ab, teste A/B, Modo A (TC) x Modo B (SE)
  scripts/generate_dataset_paper.py   gerador da camada G (semente 42)
  scripts/gerar_tabelas_texto.py      tabelas e macros (Ap. A, auditoria, 310)
  scripts/empacotar_colab.py          pacote dos cadernos, com os hashes
  scripts/lote_validacao/             coleta, alvos, redacao, anotacao e montagem
  scripts/v3/                         ciclos, autoteste do validador, decomposicao
  docs/manual_de_anotacao.md          manual de anotacao
  docs/lote_validacao.md              metodologia dos lotes
  docs/v3.md, docs/v3_desenvolvimento.md   desenho, plano e ciclos da v3
  docs/como_estender.md               guia para acrescentar uma versao
  data/dataset.json                   dataset (310 consultas)
  data/lote_validacao.json            lote 1
  data/lote_validacao_2.json          lote 2 (teste da v3)
  data/auditoria/                     artefatos da auditoria
  db/                                 esquema do prototipo e dados sinteticos
  results/                            registros e manifestos das rodadas
  results/ab/                         teste A/B (ab.md e ab.json)
  notebooks/avaliacao_colab.ipynb     rodadas completas em GPU de nuvem
  notebooks/linha_de_base_colab.ipynb linhas de base em GPU de nuvem
  notebooks/lote_validacao_colab.ipynb   rodadas do lote 1 (v1 e v2)
  notebooks/lote2_v3_colab.ipynb      rodadas do lote 2 (confere o hash da v3)
  notebooks/lote2_v3_modelos_colab.ipynb   lote 2 com Gemma 4 E2B e Qwen 3 4B
  experimentos/orientacoes/           variante descartada da v3
  tests/                              testes automatizados (pytest)
  Dockerfile, docker-compose.yml      imagem da API; servicos postgis e api
  .github/workflows/testes.yml        integracao continua (ruff, pytest, imagem)
  texto/                              fontes LaTeX e PDF deste trabalho
\end{verbatim}
}

Os passos abaixo, com os comandos descritos na Seção~\ref{apx:comandos}, reproduzem os resultados num ambiente com Python 3.12 ou superior, o Ollama e, para o banco da demonstração, o Docker.

{\footnotesize
\begin{verbatim}
git clone https://github.com/vhenequim/bdgex-busca-tool-calling.git pfc_busca
cd pfc_busca
pip install -e ".[dev]"     # pacote, comandos pfc e pfc-*, pytest e ruff
pytest -q tests/            # sem Ollama, sem Groq e sem banco
pfc dataset                 # reconstroi data/dataset.json (mesmo hash)
pfc auditar --anotacao-independente data/auditoria/anotacao_independente.json
ollama pull qwen3:4b-instruct-2507-q4_K_M
pfc avaliar --modelo qwen3:4b-instruct-2507-q4_K_M --repeticoes 3 \
            --hoje 2026-09-14 --sem-sql
pfc conferir                # rodadas das 310, da estacao e do lote 1
pfc conferir results/lote2/*-* results/v2_310/* results/v3_310/*
# tabelas e macros do texto, relidas de results/ (sem nova rodada)
python scripts/gerar_tabelas_texto.py --texto texto
pfc lote  --paper texto                             # lote 1 e v2
pfc lote2 --paper texto                             # lote 2 e v3
pfc lote2 --modelo gemma4:e2b-it-qat --paper texto  # extensao da v3
pfc lote2 --modelo qwen3:4b-instruct-2507-q4_K_M --paper texto
pfc ab --estrito --paper texto                      # teste A/B
python scripts/v3/validador_no_gabarito.py --tex    # autoteste do validador
# API com a abordagem recomendada, sobre os dados sinteticos
docker compose up -d postgis
pfc api
\end{verbatim}
}

Os comandos de análise releem os registros de execução e reaplicam o gabarito vigente, lido de \texttt{data/dataset.json} ou do arquivo do lote. Assim, as tabelas e os valores dos Capítulos~\ref{cap:resultados} e~\ref{cap:conclusao} podem ser regenerados sem nova execução dos modelos. Num clone, o autoteste do validador grava as suas macros só em \texttt{data/\allowbreak v3/\allowbreak numeros\_v3.tex}, que deve ser copiado à mão para \texttt{texto/\allowbreak tabelas/}.

O índice de folhas do BDGEx que as ferramentas da v3 consultam é gerado por \texttt{scripts/v3/gerar\_dados\_ferramentas.py} a partir da coleta do catálogo. Nem o índice nem a coleta são versionados, porque a política do BDGEx não trata de redistribuição. O índice entra no \textit{hash} da configuração da v3 gravado em cada manifesto. Uma coleta de outra data, com o catálogo já alterado, gera outro índice e outro \textit{hash}, e a v3 deixa de coincidir com a configuração avaliada. O \texttt{pfc conferir} acusa a diferença, e o \texttt{/api/health} não encontra nenhuma rodada avaliada com o mesmo \textit{hash}. Sem o índice, a API sobe em modo degradado, com aviso no \textit{log}, no \texttt{/api/health} e em cada resposta, e as ferramentas deixam de reconhecer nomes de folha e de conferir códigos no acervo. A reprodução exata da v3 exige o índice da coleta de 28/09/2026, cujas páginas têm o \textit{hash} SHA-256 registrado em \texttt{data/\allowbreak lote\_validacao/\allowbreak bdgex\_resumo.json}. A repontuação a partir dos registros não depende dele.

O repositório guarda também os relatórios que detalham o Capítulo~\ref{cap:resultados}. O \texttt{results/\allowbreak consolidado/\allowbreak erros\_representativos.md} traz até quatro erros sorteados de cada tipo por modelo, nas rodadas completas. O \texttt{results/\allowbreak ab/\allowbreak ab.md} lista, para cada par do teste A/B, as pastas das duas rodadas, as consultas e repetições e as condições de \textit{hardware} que decidem se a latência é comparável. O \texttt{results/\allowbreak lote2/\allowbreak consolidado/\allowbreak lote2.md} traz as medidas do plano do lote 2 que não estão nas tabelas, como os efeitos de forma. A pasta \texttt{results/\allowbreak \_descartados/\allowbreak estacao\_14set\_vram\_disputada/} guarda a rodada do Qwen 3 4B feita na estação em 14 de setembro de 2026 e depois substituída, que a Tabela~\ref{tab:concordancia-reexecucao} compara com a rodada final. O manifesto dessa rodada, anterior ao registro da VRAM livre, não informa a parcela do modelo na GPU, que pode ter sido outra na rodada de 24 de setembro.

Os artefatos do lote 1 (validação; Seções~\ref{sec:lote-validacao} e~\ref{sec:lote-resultados}) ficam em \texttt{data/\allowbreak lote\_validacao/}, com a metodologia em \texttt{docs/lote\_validacao.md}. Eles incluem os alvos, a redação e a anotação, as consultas descartadas com o motivo e a auditoria dos erros. O \texttt{pfc lote} calcula também, sem chamada nova ao modelo, as configurações simuladas da Tabela~\ref{tab:lote_duas_etapas}. Os artefatos do lote 2 (teste; Seção~\ref{sec:lote-validacao}) ficam em \texttt{data/\allowbreak lote\_validacao\_2/}, e as rodadas e o relatório do teste da v3, em \texttt{results/lote2/}. Os ciclos de desenvolvimento da v3 estão em \texttt{docs/v3\_desenvolvimento.md} e \texttt{results/dev\_v3/}.

\section{Comandos do \textit{pipeline} de avaliação}
\label{apx:comandos}

Esta seção completa as Seções~\ref{sec:pipeline-avaliacao} e~\ref{sec:harness}. Os comandos \texttt{pfc-dataset}, \texttt{pfc-auditar}, \texttt{pfc-avaliar}, \texttt{pfc-relatorio}, \texttt{pfc-conferir} e \texttt{pfc-ab} são pontos de entrada do pacote. O \texttt{pfc} os reúne em subcomandos com os mesmos argumentos e acrescenta \texttt{pfc abordagens}, que lista o registro e os \textit{hashes}, \texttt{pfc api}, que sobe a API, e \texttt{pfc lote} e \texttt{pfc lote2}, que analisam os lotes.

As checagens do \texttt{pfc-auditar} cobrem estrutura, rastreabilidade, evidência pelo anotador por regras e consistência. O \texttt{pfc-avaliar} grava os registros em JSON Lines. A opção \texttt{-{}-abordagem} escolhe a configuração (\texttt{saida\_estruturada}, \texttt{prototipo}, \texttt{tool\_calling\_v2}, \texttt{saida\_estruturada\_v2}, \texttt{tool\_calling\_v3d}, \texttt{tool\_calling\_v3a}, \texttt{tool\_calling\_v3} e \texttt{saida\_estruturada\_v3}), e o padrão é o \textit{Tool Calling} v1 (\texttt{tool\_calling}). A opção \texttt{-{}-dataset} aponta para o arquivo de cada lote, e \texttt{-{}-provedor groq} troca o \texttt{ChatOllama} pelo \texttt{ChatOpenAI} apontado para o \textit{endpoint} compatível do Groq (Seção~\ref{apx:specs-groq}). O \texttt{pfc-relatorio} extrai também uma amostra de erros representativos por tipo.

Na API, as variáveis \texttt{PFC\_ABORDAGEM} e \texttt{PFC\_MODELO\_PADRAO} escolhem a abordagem e o modelo, e \texttt{PFC\_ABORDAGEM=tool\_calling} volta à v1. Um nome fora do registro impede a subida. Com a v3, a SQL é executada uma vez, com os parâmetros da busca que encerra o laço de chamadas. O Docker Compose sobe a API com o banco, com o Ollama no \textit{host}.

A integração contínua (\texttt{.github/\allowbreak workflows/\allowbreak testes.yml}) roda o \texttt{ruff} e o \texttt{pytest} em Python 3.12 a cada \textit{push} e \textit{pull request} e constrói a imagem da API, que precisa subir e responder o \texttt{/api/health} sem Ollama e sem banco. Os testes cobrem também o tempo relativo, as invariantes do \textit{dataset} e os relatórios, e a API é testada com tradutores falsos.

O guia \texttt{docs/como\_estender.md} descreve, para a próxima versão, os passos que a v3 seguiu (módulo novo, entrada no registro, congelamento pelo \textit{hash}, plano de análise gravado antes da rodada e números levados ao texto por macros geradas).

% ----------------------------------------------------------
% Apêndice D — Resultados Complementares
% Tabelas e figuras geradas (mesmos comandos do cabeçalho de cap05-resultados.tex).
% ----------------------------------------------------------
\chapter{Resultados Complementares}
\label{apx:resultados}

Este apêndice reúne as tabelas e as figuras que detalham o Capítulo~\ref{cap:resultados}.

\section{Rodadas completas das 310 consultas}
\label{apx:res-310}

Os quatro modelos locais executaram as 310 consultas na T4, com três repetições (Seção~\ref{sec:comparativo}). Na Figura~\ref{fig:fig_acuracia_por_categoria}, a acurácia da categoria F é a taxa de recusa correta, porque o gabarito dessas consultas não tem campos.

\tabelagerada{tab_latencia}{latência por modelo}

\figuragerada{fig_latencia_boxplot}{Distribuição da latência da tradução nas rodadas completas, em escala logarítmica, por modelo}{boxplot de latência}

\tabelagerada{tab_mcnemar}{McNemar entre pares de modelos}

\tabelagerada{tab_por_categoria}{F1 por categoria}

\figuragerada{fig_acuracia_por_categoria}{Acurácia por consulta, por categoria e por modelo, nas rodadas completas dos modelos locais}{acurácia por categoria}

\tabelagerada{tab_por_origem}{acurácia por camada}

\tabelagerada{tab_por_campo}{F1 por campo}

\section{Estação de referência}
\label{apx:res-estacao}

Na validação na estação de referência, o Qwen 3 4B e o Gemma 4 E2B executaram as 62 consultas das camadas P e N (Seção~\ref{sec:validacao-local}). O teste de McNemar entre os dois está na Tabela~\ref{tab:mcnemar_estacao}. Na estação, a comparação entre as abordagens confirmou o empate no Gemma 4 E2B. Nessa amostra menor, ela não mostrou diferença significativa no Qwen 3 4B nem no Gemma 4 E4B, dois modelos em que a vantagem da Saída Estruturada na T4 era significativa (Tabela~\ref{tab:abordagens_estacao}).

\tabelagerada{tab_comparativo_estacao}{comparativo na estação}

\tabelagerada{tab_latencia_estacao}{latência na estação}

\tabelagerada{tab_mcnemar_estacao}{McNemar na estação}

\tabelagerada{tab_abordagens_estacao}{Tool Calling e Saída Estruturada na estação}

\tabelagerada{tab_concordancia_ambientes}{concordância entre a estação e as rodadas completas}

\tabelagerada{tab_concordancia_reexecucao}{reexecução do Qwen 3 4B na estação}

\section{Referência em nuvem}
\label{apx:res-nuvem}

Os \res{groq}{geral}{nmodelos} modelos servidos pelo Groq executaram as 310 consultas uma vez (Seção~\ref{sec:referencia-nuvem}).

\tabelagerada{tab_comparativo_groq}{comparativo da referência em nuvem}

\tabelagerada{tab_por_categoria_groq}{F1 por categoria na referência em nuvem}

\section{Lote 1 (validação)}
\label{apx:lote}

As Tabelas~\ref{tab:lote_pares}, \ref{tab:lote_duas_etapas}, \ref{tab:lote_categorias} e~\ref{tab:lote_campos} detalham a etapa 3 no lote 1, com o Gemma 4 E4B (Seção~\ref{sec:lote-resultados}).

\tabelagerada{tab_lote_pares}{Comparações pareadas no lote}

\tabelagerada{tab_lote_duas_etapas}{duas etapas e regra do objeto vazio}

\tabelagerada{tab_lote_categorias}{Acurácia por categoria no lote}

\tabelagerada{tab_lote_campos}{F1 por campo no lote}

\paragraph{Registros e leituras múltiplas.} Nas consultas do domínio, a Saída Estruturada v1 foi a melhor configuração nos oito registros de linguagem. As vantagens aparentes do \textit{Tool Calling} em registros como ``direto'' e ``sem acento'' vêm da concentração das consultas fora do domínio nesses registros. Os registros sem acento e telegráfico ficaram abaixo dos demais nas quatro configurações (na Saída Estruturada v1, \res{lote}{se1}{sematel} contra \res{lote}{se1}{sematelresto}). Todas as configurações acertaram menos nas consultas com mais de uma leitura aceita, e nelas a vantagem da Saída Estruturada cresceu (\res{lote}{multiplas}{didv1} p.p. na v1 e \res{lote}{multiplas}{didv2} p.p. na v2, na diferença entre as diferenças). A ambiguidade em si não explica esse aumento. As leituras múltiplas se concentram em algumas famílias, e, comparadas as abordagens dentro de cada família, a diferença desaparece na v2. Na v1, ela vem sobretudo da não busca do \textit{Tool Calling}, mais frequente nessas consultas.
